\input{preamble}

% OK, start here.
%
\begin{document}

\title{Discriminants et différentes}


\maketitle

\phantomsection
\label{section-phantom}

\tableofcontents

\section{Introduction}
\label{section-introduction}

\noindent
Dans ce chapitre, nous étudions la différente et le discriminant
des morphismes localement quasi-finis de schémas.
On pourra consulter \cite{Kunz} pour une partie de ces résultats.

\medskip\noindent
Étant donné un morphisme quasi-fini $f : Y \to X$ de schémas noethériens,
on dispose d'un module dualisant relatif $\omega_{Y/X}$.
Dans la section \ref{section-quasi-finite-dualizing},
nous construisons ce module directement, à l'aide
du théorème principal de Zariski et de méthodes de localisation étale.
Sa propriété essentielle est la suivante : étant donné un diagramme
$$
\xymatrix{
Y' \ar[d]_{f'} \ar[r]_{g'} & Y \ar[d]^f \\
X' \ar[r]^g & X
}
$$
où $g : X' \to X$ est plat, $Y' \subset X' \times_X Y$ est ouvert et
$f' : Y' \to X'$ est fini, il existe un isomorphisme canonique
$$
f'_*(g')^*\omega_{Y/X} =
\SheafHom_{\mathcal{O}_{X'}}(f'_*\mathcal{O}_{Y'}, \mathcal{O}_{X'})
$$
de faisceaux de $f'_*\mathcal{O}_{Y'}$-modules. Dans la
section \ref{section-quasi-finite-traces}, nous montrons que,
si $f$ est plat, il existe une section globale canonique
$\tau_{Y/X} \in H^0(Y, \omega_{Y/X})$ telle que, pour tout diagramme
commutatif comme ci-dessus, l'isomorphisme envoie $(g')^*\tau_{Y/X}$ sur l'application trace
de la section \ref{section-discriminant}
associée au morphisme fini localement libre $f'$.
Dans la section \ref{section-different},
nous définissons la différente d'un morphisme plat quasi-fini
de schémas noethériens comme l'annulateur du
conoyau de $\tau_{Y/X} : \mathcal{O}_X \to \omega_{Y/X}$.

\medskip\noindent
L'objectif principal de ce chapitre est de montrer que, pour
$f$ quasi-fini syntomique\footnote{Autrement dit plat et localement d'intersection complète.}, la
différente coïncide avec la différente de Kähler.
La différente de Kähler est l'idéal de Fitting d'ordre zéro de $\Omega_{Y/X}$, voir
la section \ref{section-kahler-different}.
Cette égalité n'est pas évidente ; nous utilisons un argument élégant
dû à Tate, voir la section \ref{section-formula-different}.
Nous étudions également, chemin faisant, la différente de Noether
et la différente de Dedekind.

\medskip\noindent
Ce n'est qu'à la fin de ce chapitre, voir
les sections \ref{section-comparison} et \ref{section-gorenstein-lci},
que nous établissons le lien avec les résultats plus avancés
sur la dualité des schémas.








\section{Modules dualisants des homomorphismes quasi-finis d'anneaux}
\label{section-quasi-finite-dualizing}

\noindent
Soit $A \to B$ un homomorphisme quasi-fini d'anneaux noethériens. D'après
le théorème principal de Zariski
(Algèbre, lemme \ref{algebra-lemma-quasi-finite-open-integral-closure}),
il existe une factorisation $A \to B' \to B$ où
$A \to B'$ est fini et $B' \to B$ induit une immersion ouverte des spectres.
On pose
\begin{equation}
\label{equation-dualizing}
\omega_{B/A} = \Hom_A(B', A) \otimes_{B'} B
\end{equation}
dans cette situation. On peut considérer ce module comme une sorte de module
dualisant relatif, voir les lemmes \ref{lemma-compare-dualizing} et
\ref{lemma-compare-dualizing-algebraic}.
Dans cette section, nous montrerons par des méthodes élémentaires d'algèbre commutative
que $\omega_{B/A}$ est indépendant du choix de la factorisation
et que la formation de $\omega_{B/A}$ commute au changement de base plat.
Pour établir l'indépendance à l'égard de la factorisation, nous comparons deux
factorisations données.

\begin{lemma}
\label{lemma-dominate-factorizations}
Soit $A \to B$ un homomorphisme quasi-fini d'anneaux. Étant donné deux factorisations
$A \to B' \to B$ et $A \to B'' \to B$ où
$A \to B'$ et $A \to B''$ sont finis et $\Spec(B) \to \Spec(B')$
et $\Spec(B) \to \Spec(B'')$ sont des immersions ouvertes, il existe
une $A$-sous-algèbre $B''' \subset B$ finie sur $A$ telle que
$\Spec(B) \to \Spec(B''')$ soit une immersion ouverte et que $B' \to B$ et
$B'' \to B$ se factorisent par $B'''$.
\end{lemma}

\begin{proof}
Soit $B''' \subset B$ la $A$-sous-algèbre engendrée par les images
de $B' \to B$ et $B'' \to B$. Comme $B'$ et $B''$ sont chacun engendrés
par un nombre fini d'éléments entiers sur $A$, on voit que $B'''$ est
engendrée par un nombre fini d'éléments entiers sur $A$, et l'on en déduit
que $B'''$ est finie sur $A$
(Algèbre, lemme \ref{algebra-lemma-characterize-finite-in-terms-of-integral}).
Considérons les applications
$$
B = B' \otimes_{B'} B \to B''' \otimes_{B'} B \to B \otimes_{B'} B = B
$$
La dernière égalité tient au fait que $\Spec(B) \to \Spec(B')$ est une
immersion ouverte (donc un monomorphisme). La seconde flèche est injective
puisque $B' \to B$ est plat. Les deux flèches sont donc des isomorphismes.
Cela signifie que
$$
\xymatrix{
\Spec(B''') \ar[d] & \Spec(B) \ar[d] \ar[l] \\
\Spec(B') & \Spec(B) \ar[l]
}
$$
est cartésien. Comme le changement de base d'une immersion ouverte est une
immersion ouverte, on conclut.
\end{proof}

\begin{lemma}
\label{lemma-dualizing-well-defined}
Le module (\ref{equation-dualizing}) est bien défini, c'est-à-dire
indépendant du choix de la factorisation.
\end{lemma}

\begin{proof}
Soient $B', B'', B'''$ comme dans le lemme \ref{lemma-dominate-factorizations}.
On obtient une application canonique
$$
\omega''' = \Hom_A(B''', A) \otimes_{B'''} B \longrightarrow
\Hom_A(B', A) \otimes_{B'} B = \omega'
$$
et une application analogue faisant intervenir $B''$. Il suffit de montrer que ces applications
sont des isomorphismes. Soit $g \in B'$ un élément tel que
$B'_g \to B_g$ soit un isomorphisme, de sorte que $B'_g \to (B''')_g \to B_g$
soient des isomorphismes. Il suffit de montrer que $(\omega''')_g \to \omega'_g$
est un isomorphisme. Le noyau et le conoyau de l'homomorphisme d'anneaux $B' \to B'''$
sont des $A$-modules de type fini et de torsion par rapport aux puissances de $g$.
Ils sont donc annulés par une puissance de $g$.
Le résultat s'en déduit aisément.
\end{proof}

\begin{lemma}
\label{lemma-localize-dualizing}
Soit $A \to B$ un homomorphisme quasi-fini d'anneaux noethériens.
\begin{enumerate}
\item Si $A \to B$ se factorise en $A \to A_f \to B$ pour un $f \in A$,
alors $\omega_{B/A} = \omega_{B/A_f}$.
\item Si $g \in B$, alors $(\omega_{B/A})_g = \omega_{B_g/A}$.
\item Si $f \in A$, alors $\omega_{B_f/A_f} = (\omega_{B/A})_f$.
\end{enumerate}
\end{lemma}

\begin{proof}
Soit $A \to B' \to B$ une factorisation où $A \to B'$ est fini et
$\Spec(B) \to \Spec(B')$ est une immersion ouverte. Dans le cas (1), on peut utiliser
la factorisation $A_f \to B'_f \to B$ pour calculer $\omega_{B/A_f}$
et appliquer Algèbre, lemme \ref{algebra-lemma-hom-from-finitely-presented}.
Dans le cas (2), on utilise la factorisation $A \to B' \to B_g$ pour obtenir le résultat.
L'assertion (3) résulte de (1) et (2).
\end{proof}

\noindent
Soit $A \to B$ un homomorphisme quasi-fini d'anneaux noethériens, soit
$A \to A_1$ un homomorphisme quelconque d'anneaux noethériens, et posons
$B_1 = B \otimes_A A_1$. On obtient un diagramme cocartésien
$$
\xymatrix{
B \ar[r] & B_1 \\
A \ar[u] \ar[r] & A_1 \ar[u]
}
$$
Remarquons que $A_1 \to B_1$ est également quasi-fini (Algèbre, lemme
\ref{algebra-lemma-quasi-finite-base-change}).
Dans cette situation, nous allons définir une application canonique
$B$-linéaire de changement de base
\begin{equation}
\label{equation-bc-dualizing}
\omega_{B/A} \longrightarrow \omega_{B_1/A_1}
\end{equation}
À cet effet, on choisit une factorisation $A \to B' \to B$ comme dans la construction
de $\omega_{B/A}$. Alors $B'_1 = B' \otimes_A A_1$ est fini sur $A_1$
et l'on peut utiliser la factorisation $A_1 \to B'_1 \to B_1$ pour construire
$\omega_{B_1/A_1}$. Il s'agit donc de construire une application
$$
\Hom_A(B', A) \otimes_{B'} B
\longrightarrow
\Hom_{A_1}(B' \otimes_A A_1, A_1) \otimes_{B'_1} B_1
$$
Il suffit ainsi de construire une application $B'$-linéaire\newline
$\Hom_A(B', A) \to \Hom_{A_1}(B' \otimes_A A_1, A_1)$,
que l'on notera $\varphi \mapsto \varphi_1$.
Étant donné une application $A$-linéaire $\varphi : B' \to A$,
on prend pour $\varphi_1$ l'application définie par
$\varphi_1(b' \otimes a_1) = \varphi(b')a_1$.
Elle est manifestement $A_1$-linéaire, ce qui achève la construction.

\begin{lemma}
\label{lemma-bc-map-dualizing}
L'application de changement de base (\ref{equation-bc-dualizing})
est indépendante du choix de la
factorisation $A \to B' \to B$. Étant donné des homomorphismes d'anneaux $A \to A_1 \to A_2$,
la composée des applications de changement de base pour $A \to A_1$ et $A_1 \to A_2$
est l'application de changement de base pour $A \to A_2$.
\end{lemma}

\begin{proof}
Omise. Indication : raisonner exactement comme dans le
lemme \ref{lemma-dualizing-well-defined},
en utilisant le lemme \ref{lemma-dominate-factorizations}.
\end{proof}

\begin{lemma}
\label{lemma-dualizing-flat-base-change}
Si $A \to A_1$ est plat, alors
l'application de changement de base (\ref{equation-bc-dualizing}) induit un isomorphisme
$\omega_{B/A} \otimes_B B_1 \to \omega_{B_1/A_1}$.
\end{lemma}

\begin{proof}
Supposons $A \to A_1$ plat. Par construction de $\omega_{B/A}$, on peut
supposer que $A \to B$ est fini. Alors $\omega_{B/A} = \Hom_A(B, A)$ et
$\omega_{B_1/A_1} = \Hom_{A_1}(B_1, A_1)$. Puisque $B_1 = B \otimes_A A_1$,
le résultat découle de Compléments d'algèbre, lemme
\ref{more-algebra-lemma-pseudo-coherence-and-base-change-ext}.
\end{proof}

\begin{lemma}
\label{lemma-dualizing-composition}
Soient $A \to B \to C$ des homomorphismes quasi-finis d'anneaux noethériens.
Il existe une application canonique
$\omega_{B/A} \otimes_B \omega_{C/B} \to \omega_{C/A}$.
\end{lemma}

\begin{proof}
Choisissons $A \to B' \to B$ où $A \to B'$ est fini et
$\Spec(B) \to \Spec(B')$ est une immersion ouverte. Alors
$B' \to C$ est également quasi-fini. Choisissons $B' \to C' \to C$
où $B' \to C'$ est fini et $\Spec(C) \to \Spec(C')$ est une
immersion ouverte. La source de la flèche est alors
$$
\Hom_A(B', A) \otimes_{B'} B \otimes_B
\Hom_B(B \otimes_{B'} C', B) \otimes_{B \otimes_{B'} C'} C
$$
qui est égale à
$$
\Hom_A(B', A) \otimes_{B'}
\Hom_{B'}(C', B) \otimes_{C'} C
$$
Ce module est bien muni d'une application canonique vers
$\Hom_A(C', A) \otimes_{C'} C = \omega_{C/A}$,
provenant de la composition
$\Hom_A(B', A) \times \Hom_{B'}(C', B) \to \Hom_A(C', A)$.
\end{proof}

\begin{lemma}
\label{lemma-dualizing-product}
Soient $A \to B$ et $A \to C$ des homomorphismes quasi-finis d'anneaux noethériens.
Alors $\omega_{B \times C/A} = \omega_{B/A} \times \omega_{C/A}$
comme modules sur $B \times C$.
\end{lemma}

\begin{proof}
Choisissons des factorisations $A \to B' \to B$ et $A \to C' \to C$ telles que
$A \to B'$ et $A \to C'$ soient finis et que $\Spec(B) \to \Spec(B')$
et $\Spec(C) \to \Spec(C')$ soient des immersions ouvertes. Alors
$A \to B' \times C' \to B \times C$ est une factorisation analogue.
Le calcul de $\omega_{B \times C/A}$ à l'aide de cette factorisation
donne le lemme.
\end{proof}

\begin{lemma}
\label{lemma-dualizing-associated-primes}
Soit $A \to B$ un homomorphisme quasi-fini d'anneaux noethériens.
Alors $\text{Ass}_B(\omega_{B/A})$ est l'ensemble des idéaux premiers de $B$
au-dessus d'idéaux premiers associés de $A$.
\end{lemma}

\begin{proof}
Choisissons une factorisation $A \to B' \to B$ où $A \to B'$ est fini et
$B' \to B$ induit une immersion ouverte des spectres. Comme
$\omega_{B/A} = \omega_{B'/A} \otimes_{B'} B$, il suffit
de démontrer l'assertion pour $\omega_{B'/A}$. On peut donc supposer $A \to B$
fini.

\medskip\noindent
Supposons $\mathfrak p \in \text{Ass}(A)$ et soit $\mathfrak q$ un idéal premier
de $B$ au-dessus de $\mathfrak p$. Soit $x \in A$ un élément dont
l'annulateur est $\mathfrak p$. Choisissons une application $\kappa(\mathfrak p)$-linéaire
non nulle $\lambda : \kappa(\mathfrak q) \to \kappa(\mathfrak p)$.
Puisque $A/\mathfrak p \subset B/\mathfrak q$ est une extension finie
d'anneaux, il existe $f \in A$, $f \not \in \mathfrak p$,
tel que $f\lambda$ envoie $B/\mathfrak q$ dans $A/\mathfrak p$.
On obtient donc une application $A$-linéaire non nulle
$$
B \to B/\mathfrak q \to A/\mathfrak p \to A,\quad
b \mapsto f\lambda(b)x
$$
Un calcul immédiat montre que cet élément de $\omega_{B/A}$
a pour annulateur $\mathfrak q$, d'où
$\mathfrak q \in \text{Ass}(\omega_{B/A})$.

\medskip\noindent
Réciproquement, supposons que $\mathfrak q \subset B$ soit un idéal premier
au-dessus d'un idéal premier $\mathfrak p \subset A$ qui n'est pas associé
à $A$. Il faut montrer que
$\mathfrak q \not \in \text{Ass}_B(\omega_{B/A})$.
En remplaçant $A$ par $A_\mathfrak p$ et $B$ par
$B_\mathfrak p$, on peut supposer que $\mathfrak p$ est un idéal maximal
de $A$. Ceci est permis par le lemme \ref{lemma-dualizing-flat-base-change} et
Algèbre, lemme \ref{algebra-lemma-localize-ass}.
Il existe alors $f \in \mathfrak m$
qui est non diviseur de zéro dans $A$.
Cet élément $f$ est non diviseur de zéro sur $\omega_{B/A}$,
et $\mathfrak q$ n'est donc pas un idéal premier associé à ce module.
\end{proof}

\begin{lemma}
\label{lemma-dualizing-base-flat-flat}
Soit $A \to B$ un homomorphisme plat quasi-fini d'anneaux noethériens.
Alors $\omega_{B/A}$ est un $A$-module plat.
\end{lemma}

\begin{proof}
Soit $\mathfrak q \subset B$ un idéal premier au-dessus de $\mathfrak p \subset A$.
Nous allons montrer que le localisé $\omega_{B/A, \mathfrak q}$ est plat
sur $A_\mathfrak p$.
Cela suffit d'après Algèbre, lemme \ref{algebra-lemma-flat-localization}.
D'après
Algèbre, lemme \ref{algebra-lemma-etale-makes-quasi-finite-finite-one-prime},
on peut trouver un homomorphisme étale d'anneaux $A \to A'$ et un idéal
premier $\mathfrak p' \subset A'$ au-dessus de $\mathfrak p$
tels que $\kappa(\mathfrak p') = \kappa(\mathfrak p)$ et
que
$$
B' = B \otimes_A A' = C \times D
$$
avec $A' \to C$ fini, et tels que l'unique idéal premier $\mathfrak q'$
de $B \otimes_A A'$ au-dessus de $\mathfrak q$ et de $\mathfrak p'$
corresponde à un idéal premier de $C$. D'après le
lemme \ref{lemma-dualizing-flat-base-change}
et Algèbre, lemme \ref{algebra-lemma-base-change-flat-up-down},
il suffit de montrer que $\omega_{B'/A', \mathfrak q'}$
est plat sur $A'_{\mathfrak p'}$.
Puisque $\omega_{B'/A'} = \omega_{C/A'} \times \omega_{D/A'}$
d'après le lemme \ref{lemma-dualizing-product},
on est ramené au cas où $B$ est fini plat sur $A$.
Dans ce cas, $B$ est un $A$-module fini localement libre
et $\omega_{B/A} = \Hom_A(B, A)$ est le $A$-module dual,
fini localement libre.
\end{proof}

\begin{lemma}
\label{lemma-dualizing-base-change-of-flat}
Si $A \to B$ est plat, l'application de changement de base (\ref{equation-bc-dualizing})
induit un isomorphisme $\omega_{B/A} \otimes_B B_1 \to \omega_{B_1/A_1}$.
\end{lemma}

\begin{proof}
Si $A \to B$ est fini plat, $B$ est un $A$-module fini localement libre.
Dans ce cas, $\omega_{B/A} = \Hom_A(B, A)$ est le $A$-module dual,
fini localement libre, et sa formation commute
à tout changement de base, ce qui démontre le lemme dans ce cas.
Dans le paragraphe suivant, nous ramenons le cas général (quasi-fini plat)
au cas fini plat que l'on vient de traiter.

\medskip\noindent
Soit $\mathfrak q_1 \subset B_1$ un idéal premier. Nous allons montrer que la
localisation de l'application en $\mathfrak q_1$ est un isomorphisme, ce qui
suffit d'après Algèbre, lemme \ref{algebra-lemma-characterize-zero-local}.
Soient $\mathfrak q \subset B$ et $\mathfrak p \subset A$ les idéaux
premiers au-dessous de $\mathfrak q_1$. D'après
Algèbre, lemme \ref{algebra-lemma-etale-makes-quasi-finite-finite-one-prime},
on peut trouver un homomorphisme étale d'anneaux $A \to A'$ et un idéal
premier $\mathfrak p' \subset A'$ au-dessus de $\mathfrak p$
tels que $\kappa(\mathfrak p') = \kappa(\mathfrak p)$ et
que
$$
B' = B \otimes_A A' = C \times D
$$
avec $A' \to C$ fini, et tels que l'unique idéal premier $\mathfrak q'$
de $B \otimes_A A'$ au-dessus de $\mathfrak q$ et de $\mathfrak p'$
corresponde à un idéal premier de $C$. Posons $A'_1 = A' \otimes_A A_1$ et
considérons les applications de changement de base
(\ref{equation-bc-dualizing}) pour les homomorphismes d'anneaux
$A \to A' \to A'_1$ et $A \to A_1 \to A'_1$, dans le diagramme
$$
\xymatrix{
\omega_{B'/A'} \otimes_{B'} B'_1 \ar[r] & \omega_{B'_1/A'_1} \\
\omega_{B/A} \otimes_B B'_1 \ar[r] \ar[u] &
\omega_{B_1/A_1} \otimes_{B_1} B'_1 \ar[u]
}
$$
où $B' = B \otimes_A A'$, $B_1 = B \otimes_A A_1$, et
$B_1' = B \otimes_A (A' \otimes_A A_1)$. D'après le
lemme \ref{lemma-bc-map-dualizing}, le diagramme commute. D'après le
lemme \ref{lemma-dualizing-flat-base-change},
les flèches verticales sont des isomorphismes.
Comme $B_1 \to B'_1$ est étale, donc plat, il suffit
de montrer que la flèche horizontale supérieure est un isomorphisme après localisation
en un idéal premier $\mathfrak q'_1$ de $B'_1$ au-dessus de $\mathfrak q$
(un tel idéal premier existe ; appliquer
Algèbre, lemme \ref{algebra-lemma-local-flat-ff}).
On peut donc supposer que $B = C \times D$, avec $A \to C$
fini et $\mathfrak q$ correspondant à un idéal premier de $C$.
Dans ce cas, le module dualisant $\omega_{B/A}$ se décompose
de façon analogue (lemme \ref{lemma-dualizing-product}),
ce qui ramène la question
au cas fini plat $A \to C$ traité ci-dessus.
\end{proof}

\begin{remark}
\label{remark-relative-dualizing-for-quasi-finite}
Soit $f : Y \to X$ un morphisme localement quasi-fini de schémas localement
noethériens. Il résulte immédiatement du lemme \ref{lemma-localize-dualizing}
qu'il existe un unique $\mathcal{O}_Y$-module cohérent
$\omega_{Y/X}$ sur $Y$ tel que, pour tout couple d'ouverts affines
$\Spec(B) = V \subset Y$, $\Spec(A) = U \subset X$ vérifiant $f(V) \subset U$,
on ait un isomorphisme canonique
$$
H^0(V, \omega_{Y/X}) = \omega_{B/A}
$$
et que ces isomorphismes soient compatibles aux applications de restriction.
\end{remark}

\begin{lemma}
\label{lemma-compare-dualizing-algebraic}
Soit $A \to B$ un homomorphisme quasi-fini d'anneaux noethériens.
Soit $\omega_{B/A}^\bullet \in D(B)$ le complexe dualisant relatif
algébrique étudié dans Complexes dualisants, section
\ref{dualizing-section-relative-dualizing-complexes-Noetherian}.
Il existe alors un isomorphisme (non unique)
$\omega_{B/A} = H^0(\omega_{B/A}^\bullet)$.
\end{lemma}

\begin{proof}
Choisissons une factorisation $A \to B' \to B$
où $A \to B'$ est fini et $\Spec(B') \to \Spec(B)$
est une immersion ouverte. Alors
$\omega_{B/A}^\bullet = \omega_{B'/A}^\bullet \otimes_B^\mathbf{L} B'$
d'après Complexes dualisants, lemmes
\ref{dualizing-lemma-composition-shriek-algebraic} et
\ref{dualizing-lemma-upper-shriek-localize}, et
la définition de $\omega_{B/A}^\bullet$. Il suffit donc
de montrer qu'il existe un isomorphisme lorsque $A \to B$ est fini.
Dans ce cas, on peut utiliser
Complexes dualisants, lemme \ref{dualizing-lemma-upper-shriek-finite},
pour voir que $\omega_{B/A}^\bullet = R\Hom(B, A)$, d'où
$H^0(\omega^\bullet_{B/A}) = \Hom_A(B, A)$, comme voulu.
\end{proof}





\section{Discriminant d'un morphisme fini localement libre}
\label{section-discriminant}

\noindent
Soient $X$ un schéma et $\mathcal{F}$ un $\mathcal{O}_X$-module fini localement
libre. Il existe alors une application canonique, la {\it trace},
$$
\text{Trace} :
\SheafHom_{\mathcal{O}_X}(\mathcal{F}, \mathcal{F})
\longrightarrow
\mathcal{O}_X
$$
Voir Exercices, exercice \ref{exercises-exercise-trace-det}. Cette application vérifie
la propriété suivante : $\text{Trace}(\text{id})$ est la fonction localement constante
de $\mathcal{O}_X$ correspondant au rang de $\mathcal{F}$.

\medskip\noindent
Soit $\pi : X \to Y$ un morphisme fini localement
libre de schémas. Il existe une {\it trace de $\pi$} canonique,
qui est une application $\mathcal{O}_Y$-linéaire
$$
\text{Trace}_\pi : \pi_*\mathcal{O}_X \longrightarrow \mathcal{O}_Y
$$
envoyant une section locale $f$ de $\pi_*\mathcal{O}_X$ sur la
trace de la multiplication par $f$ sur $\pi_*\mathcal{O}_X$. Sur les
ouverts affines, on retrouve la construction de
Exercices, exercice \ref{exercises-exercise-trace-det-rings}.
La composée
$$
\mathcal{O}_Y \xrightarrow{\pi^\sharp} \pi_*\mathcal{O}_X
\xrightarrow{\text{Trace}_\pi} \mathcal{O}_Y
$$
est la multiplication par le degré de $\pi$ (qui est une fonction localement constante
sur $Y$). Par analogie avec
Corps, section \ref{fields-section-trace-pairing},
on peut définir la forme trace
$$
Q_\pi :
\pi_*\mathcal{O}_X \times \pi_*\mathcal{O}_X
\longrightarrow
\mathcal{O}_Y
$$
par $(f, g) \mapsto \text{Trace}_\pi(fg)$. On peut considérer
$Q_\pi$ comme une application linéaire
$\pi_*\mathcal{O}_X \to
\SheafHom_{\mathcal{O}_Y}(\pi_*\mathcal{O}_X, \mathcal{O}_Y)$
entre modules localement libres de même rang, et l'on obtient ainsi
un déterminant
$$
\det(Q_\pi) :
\wedge^{top}(\pi_*\mathcal{O}_X)
\longrightarrow
\wedge^{top}(\pi_*\mathcal{O}_X)^{\otimes -1}
$$
ou, en d'autres termes, une section globale
$$
\det(Q_\pi) \in \Gamma(Y, \wedge^{top}(\pi_*\mathcal{O}_X)^{\otimes -2})
$$
Le {\it discriminant de $\pi$} est, par définition, le sous-schéma
fermé $D_\pi \subset Y$ défini par cette section globale.
Il est clair que $D_\pi$ est un sous-schéma fermé localement principal de $Y$.

\begin{lemma}
\label{lemma-discriminant}
Soit $\pi : X \to Y$ un morphisme fini localement
libre de schémas. Alors $\pi$ est étale si et seulement si son discriminant est vide.
\end{lemma}

\begin{proof}
D'après Morphismes, lemme \ref{morphisms-lemma-etale-flat-etale-fibres},
il suffit de vérifier que les fibres de $\pi$ sont étales.
Puisque la construction de la forme trace commute au changement de
base, on est ramené à la question suivante : soient $k$ un corps
et $A$ une $k$-algèbre de dimension finie. Montrer que
$A$ est étale sur $k$ si et seulement si la forme trace
$Q_{A/k} : A \times A \to k$, $(a, b) \mapsto \text{Trace}_{A/k}(ab)$,
est non dégénérée.

\medskip\noindent
Supposons $Q_{A/k}$ non dégénérée. Si $a \in A$ est nilpotent, alors
$ab$ est nilpotent pour tout $b \in A$, et $Q_{A/k}(a, -)$ est donc
identiquement nulle. Ainsi $A$ est réduit. On peut alors écrire
$A = K_1 \times \ldots \times K_n$ comme un produit où chaque $K_i$
est un corps (voir
Algèbre, lemmes \ref{algebra-lemma-finite-dimensional-algebra},
\ref{algebra-lemma-artinian-finite-length} et
\ref{algebra-lemma-minimal-prime-reduced-ring}).
Dans ce cas, l'espace
quadratique $(A, Q_{A/k})$ est la somme directe orthogonale des espaces
$(K_i, Q_{K_i/k})$. Il résulte de
Corps, lemme \ref{fields-lemma-separable-trace-pairing},
que chaque $K_i$ est séparable sur $k$. Cela signifie que $A$ est étale
sur $k$, d'après Algèbre, lemme \ref{algebra-lemma-etale-over-field}.
On démontre la réciproque en parcourant l'argument en sens inverse.
\end{proof}





\section{Traces des homomorphismes plats quasi-finis d'anneaux}
\label{section-quasi-finite-traces}

\noindent
La trace dont il est question dans le titre de cette section est de nature tout
autre que celle étudiée dans
Dualité des schémas, section \ref{duality-section-trace}.
Il s'agit ici de la trace
étudiée dans Corps, section \ref{fields-section-trace-pairing},
et généralisée dans Exercices, exercices \ref{exercises-exercise-trace-det} et
\ref{exercises-exercise-trace-det-rings}.

\medskip\noindent
Soit $A \to B$ un homomorphisme fini plat d'anneaux noethériens. Alors $B$ est un
$A$-module fini plat, donc fini localement libre
(Algèbre, lemme \ref{algebra-lemma-finite-projective}).
Étant donné $b \in B$, on peut considérer la {\it trace} $\text{Trace}_{B/A}(b)$
de l'application $A$-linéaire $B \to B$ donnée par
la multiplication par $b$ sur $B$. D'après les références ci-dessus, cela définit
une application $A$-linéaire $\text{Trace}_{B/A} : B \to A$.
Puisque $\omega_{B/A} = \Hom_A(B, A)$, car $A \to B$ est fini, on voit
que $\text{Trace}_{B/A} \in \omega_{B/A}$.

\medskip\noindent
Pour un homomorphisme plat quasi-fini d'anneaux quelconque, on définit la notion
de trace comme suit.

\begin{definition}
\label{definition-trace-element}
Soit $A \to B$ un homomorphisme plat quasi-fini d'anneaux noethériens.
L'{\it élément trace} est l'unique\footnote{L'unicité
et l'existence seront justifiées dans les
lemmes \ref{lemma-trace-unique} et \ref{lemma-dualizing-tau}.}
élément
$\tau_{B/A} \in \omega_{B/A}$
ayant la propriété suivante : pour toute $A$-algèbre noethérienne $A_1$
telle que $B_1 = B \otimes_A A_1$ soit muni d'une
décomposition en produit $B_1 = C \times D$ avec $A_1 \to C$ fini,
l'image de $\tau_{B/A}$ dans $\omega_{C/A_1}$
est $\text{Trace}_{C/A_1}$.
On utilise ici l'application de changement de base (\ref{equation-bc-dualizing}) et
le lemme \ref{lemma-dualizing-product} pour obtenir
$\omega_{B/A} \to \omega_{B_1/A_1} \to \omega_{C/A_1}$.
\end{definition}

\noindent
Nous montrons d'abord l'unicité des éléments trace, puis
leur existence.

\begin{lemma}
\label{lemma-trace-unique}
Soit $A \to B$ un homomorphisme plat quasi-fini d'anneaux noethériens.
Il existe alors au plus un élément trace dans $\omega_{B/A}$.
\end{lemma}

\begin{proof}
Soit $\mathfrak q \subset B$ un idéal premier au-dessus de l'idéal premier
$\mathfrak p \subset A$. D'après
Algèbre, lemme \ref{algebra-lemma-etale-makes-quasi-finite-finite-one-prime},
on peut trouver un homomorphisme étale d'anneaux $A \to A_1$ et un idéal
premier $\mathfrak p_1 \subset A_1$ au-dessus de $\mathfrak p$
tels que $\kappa(\mathfrak p_1) = \kappa(\mathfrak p)$ et
que
$$
B_1 = B \otimes_A A_1 = C \times D
$$
avec $A_1 \to C$ fini, et tels que l'unique idéal premier $\mathfrak q_1$
de $B \otimes_A A_1$ au-dessus de $\mathfrak q$ et de $\mathfrak p_1$
corresponde à un idéal premier de $C$. Remarquons que
$\omega_{C/A_1} = \omega_{B/A} \otimes_B C$
(combiner les lemmes \ref{lemma-dualizing-flat-base-change} et
\ref{lemma-dualizing-product}). Puisque la collection
des homomorphismes d'anneaux $B \to C$ ainsi obtenus forme une famille
conjointement injective d'homomorphismes plats et que l'image de $\tau_{B/A}$
dans $\omega_{C/A_1}$ est prescrite, l'unicité en résulte.
\end{proof}

\noindent
Vérifions la cohérence de cette définition.

\begin{lemma}
\label{lemma-finite-flat-trace}
Soit $A \to B$ un homomorphisme fini plat d'anneaux noethériens.
Alors $\text{Trace}_{B/A} \in \omega_{B/A}$ est l'élément trace.
\end{lemma}

\begin{proof}
Supposons donnés $A \to A_1$, avec $A_1$ noethérien, et
une décomposition en produit $B \otimes_A A_1 = C \times D$ avec $A_1 \to C$
fini. Dans ce cas, $A_1 \to D$ est évidemment fini lui aussi.
Posons $B_1 = B \otimes_A A_1$.
Puisque la construction des traces commute au changement de base,
on voit que $\text{Trace}_{B/A}$ s'envoie sur $\text{Trace}_{B_1/A_1}$.
Il suffit donc de remarquer que
$\text{Trace}_{B_1/A_1} = (\text{Trace}_{C/A_1}, \text{Trace}_{D/A_1})$
par l'isomorphisme
$\omega_{B_1/A_1} = \omega_{C/A_1} \times \omega_{D/A_1}$
du lemme \ref{lemma-dualizing-product}.
\end{proof}

\begin{lemma}
\label{lemma-trace-base-change}
Soit $A \to B$ un homomorphisme plat quasi-fini d'anneaux noethériens.
Soit $\tau \in \omega_{B/A}$ un élément trace.
\begin{enumerate}
\item Si $A \to A_1$ est un homomorphisme où $A_1$ est noethérien, alors, en posant
$B_1 = A_1 \otimes_A B$, l'image de $\tau$ dans $\omega_{B_1/A_1}$ est un
élément trace.
\item Si $A = R_f$ pour un anneau $R$ et un $f \in R$, alors
$\tau$ est un élément trace dans $\omega_{B/R}$.
\item Si $g \in B$, l'image de $\tau$ dans $\omega_{B_g/A}$
est un élément trace.
\item Si $B = B_1 \times B_2$, alors $\tau$ s'envoie sur un élément trace
dans chacun des modules $\omega_{B_1/A}$ et $\omega_{B_2/A}$.
\end{enumerate}
\end{lemma}

\begin{proof}
L'assertion (1) est une conséquence formelle de la définition.

\medskip\noindent
L'assertion (2) a un sens car $\omega_{B/R} = \omega_{B/A}$
d'après le lemme \ref{lemma-localize-dualizing}. Notons $\tau'$ l'élément
$\tau$ considéré dans $\omega_{B/R}$. Pour démontrer (2),
supposons donnés $R \to R_1$, avec $R_1$ noethérien, et une décomposition
en produit $B \otimes_R R_1 = C \times D$ où $R_1 \to C$ est fini.
En posant $A_1 = (R_1)_f$, on a $B \otimes_A A_1 = C \times D$.
Puisque $R_1 \to C$ est fini, $A_1 \to C$ l'est a fortiori.
On peut donc utiliser la propriété qui définit $\tau$ pour obtenir la propriété
correspondante de $\tau'$.

\medskip\noindent
L'assertion (3) a un sens car $\omega_{B_g/A} = (\omega_{B/A})_g$
d'après le lemme \ref{lemma-localize-dualizing}. La démonstration est analogue à celle
de (2). Supposons donnés $A \to A_1$, avec $A_1$ noethérien, et
une décomposition en produit $B_g \otimes_A A_1 = C \times D$ où $A_1 \to C$
est fini. Posons $B_1 = B \otimes_A A_1$. Alors
$\Spec(C) \to \Spec(B_1)$ est une immersion ouverte puisque $B_g \otimes_A A_1 = (B_1)_g$,
et son image est fermée puisque $B_1 \to C$ est fini
(car $A_1 \to C$ est fini).
On a donc $B_1 = C \times D_1$ et $D = (D_1)_g$. On peut alors utiliser
la propriété qui définit $\tau$ pour obtenir la propriété correspondante
de l'image de $\tau$ dans $\omega_{B_g/A}$.

\medskip\noindent
L'assertion (4) a un sens car
$\omega_{B/A} = \omega_{B_1/A} \times \omega_{B_2/A}$ d'après le
lemme \ref{lemma-dualizing-product}.
Supposons donnés $A \to A'$, avec $A'$ noethérien, et
une décomposition en produit $B \otimes_A A' = C \times D$ où $A' \to C$
est fini. On peut manifestement raffiner cette décomposition
en produit en $B \otimes_A A' = C_1 \times C_2 \times D_1 \times D_2$,
avec $A' \to C_i$ fini, de sorte que $B_i \otimes_A A' = C_i \times D_i$.
On peut alors utiliser la propriété qui définit $\tau$ pour obtenir la propriété
correspondante de l'image de $\tau$ dans $\omega_{B_i/A}$. On utilise ici le fait évident
que
$\text{Trace}_{C/A'} = (\text{Trace}_{C_1/A'}, \text{Trace}_{C_2/A'})$
pour la décomposition
$\omega_{C/A'} = \omega_{C_1/A'} \times \omega_{C_2/A'}$.
\end{proof}

\begin{lemma}
\label{lemma-glue-trace}
Soit $A \to B$ un homomorphisme plat quasi-fini d'anneaux noethériens.
Soient $g_1, \ldots, g_m \in B$ des éléments engendrant l'idéal unité.
Soit $\tau \in \omega_{B/A}$ un élément dont l'image dans
$\omega_{B_{g_i}/A}$ est un élément trace pour $A \to B_{g_i}$.
Alors $\tau$ est un élément trace.
\end{lemma}

\begin{proof}
Supposons donnés $A \to A_1$, avec $A_1$ noethérien, et une décomposition
en produit $B \otimes_A A_1 = C \times D$ où $A_1 \to C$ est fini.
Il faut montrer que l'image de $\tau$ dans $\omega_{C/A_1}$ est
$\text{Trace}_{C/A_1}$. Remarquons que $g_1, \ldots, g_m$
engendrent l'idéal unité de $B_1 = B \otimes_A A_1$ et que
$\tau$ s'envoie sur un élément trace dans $\omega_{(B_1)_{g_i}/A_1}$
d'après le lemme \ref{lemma-trace-base-change}. On peut donc remplacer
$A$ par $A_1$ et $B$ par $B_1$ pour se placer dans la situation décrite
au paragraphe suivant.

\medskip\noindent
On suppose ici que $B = C \times D$ avec $A \to C$ fini.
Soit $\tau_C$ l'image de $\tau$ dans $\omega_{C/A}$.
Il faut démontrer que $\tau_C = \text{Trace}_{C/A}$ dans $\omega_{C/A}$.
La compatibilité des éléments trace aux produits
(lemme \ref{lemma-trace-base-change})
montre que $\tau_C$ s'envoie sur un élément trace dans $\omega_{C_{g_i}/A}$.
Ainsi, en remplaçant $B$ par $C$, on peut supposer $A \to B$
fini plat.

\medskip\noindent
Supposons $A \to B$ fini plat. Dans ce cas, $\text{Trace}_{B/A}$
est un élément trace d'après le lemme \ref{lemma-finite-flat-trace}.
Par conséquent, $\text{Trace}_{B/A}$ s'envoie sur un élément trace dans
$\omega_{B_{g_i}/A}$ d'après le lemme \ref{lemma-trace-base-change}.
L'unicité des éléments trace (lemme \ref{lemma-trace-unique})
montre que $\text{Trace}_{B/A}$ et $\tau$ ont
les mêmes images dans $\omega_{B_{g_i}/A} = (\omega_{B/A})_{g_i}$.
Comme $g_1, \ldots, g_m$ engendrent l'idéal unité de $B$, l'application
$\omega_{B/A} \to \prod \omega_{B_{g_i}/A}$ est injective,
et l'on conclut que $\tau_C = \text{Trace}_{B/A}$, comme voulu.
\end{proof}

\begin{lemma}
\label{lemma-dualizing-tau}
Soit $A \to B$ un homomorphisme plat quasi-fini d'anneaux noethériens.
Il existe un élément trace $\tau \in \omega_{B/A}$.
\end{lemma}

\begin{proof}
Choisissons une factorisation $A \to B' \to B$ où $A \to B'$ est fini et
$\Spec(B) \to \Spec(B')$ est une immersion ouverte. Soient $g_1, \ldots, g_n \in B'$
des éléments tels que $\Spec(B) = \bigcup D(g_i)$ comme ouverts de $\Spec(B')$.
Supposons que l'on puisse montrer l'existence d'éléments trace $\tau_i$ pour les
homomorphismes quasi-finis plats d'anneaux $A \to B_{g_i}$. Alors, pour tous $i, j$, les éléments
$\tau_i$ et $\tau_j$ s'envoient sur des éléments trace de $\omega_{B_{g_ig_j}/A}$
d'après le lemme \ref{lemma-trace-base-change}. Par unicité des
éléments trace (lemme \ref{lemma-trace-unique}), ils ont la même image.
La condition de faisceau pour le module quasi-cohérent associé à
$\omega_{B/A}$ (voir Algèbre, lemme \ref{algebra-lemma-cover-module})
fournit donc un élément $\tau \in \omega_{B/A}$.
Alors $\tau$ est un élément trace d'après le
lemme \ref{lemma-glue-trace}.
On est ainsi ramené au cas traité au paragraphe suivant.

\medskip\noindent
Supposons $A \to B'$ fini et $g \in B'$ tel que $B = B'_g$ soit plat sur $A$.
Il s'agit de construire un élément trace dans
$\omega_{B/A} = \Hom_A(B', A) \otimes_{B'} B$.
Choisissons une résolution $F_1 \to F_0 \to B' \to 0$ de $B'$ par des
$A$-modules libres de type fini $F_0$ et $F_1$. On a alors une suite exacte
$$
0 \to \Hom_A(B', A) \to F_0^\vee \to F_1^\vee
$$
où $F_i^\vee = \Hom_A(F_i, A)$ est le module dual, libre de type fini.
De même, on a la suite exacte
$$
0 \to \Hom_A(B', B') \to F_0^\vee \otimes_A B' \to F_1^\vee \otimes_A B'
$$
L'idée de la construction de $\tau$ est d'utiliser le diagramme
$$
B' \xrightarrow{\mu} \Hom_A(B', B')
\leftarrow \Hom_A(B', A) \otimes_A B'
\xrightarrow{ev} A
$$
où la première flèche envoie $b' \in B'$ sur l'opérateur $A$-linéaire
de multiplication par $b'$, et la dernière flèche est l'application d'évaluation.
La difficulté est que la flèche du milieu, qui envoie $\lambda' \otimes b'$
sur l'application $b'' \mapsto \lambda'(b'')b'$, n'est pas un isomorphisme.
Si $B'$ est plat sur $A$, les suites exactes ci-dessus montrent que cette flèche
est un isomorphisme, et la composée de gauche à droite est la trace usuelle
$\text{Trace}_{B'/A}$. Dans le cas général, on considère
le diagramme
$$
\xymatrix{
& \Hom_A(B', A) \otimes_A B' \ar[r] \ar[d] &
\Hom_A(B', A) \otimes_A B'_g \ar[d] \\
B' \ar[r]_-\mu \ar@{..>}[rru] \ar@{..>}[ru]^\psi &
\Hom_A(B', B') \ar[r] &
\Ker(F_0^\vee \otimes_A B'_g \to F_1^\vee \otimes_A B'_g)
}
$$
La platitude de $A \to B'_g$ montre que la flèche verticale de droite est un
isomorphisme. On obtient donc la flèche pointillée sans étiquette.
Puisque $B'_g = \colim \frac{1}{g^n}B'$, que
les colimites commutent aux produits tensoriels,
et que $B'$ est un $A$-module de présentation finie,
on peut trouver $n \geq 0$ et une application $B'$-linéaire (pour la structure de $B'$-module à droite)
$\psi : B' \to \Hom_A(B', A) \otimes_A B'$
dont la composée avec la flèche verticale de gauche est $g^n\mu$.
En composant avec $ev$, on obtient un élément
$ev \circ \psi \in \Hom_A(B', A)$. On pose alors
$$
\tau = (ev \circ \psi) \otimes g^{-n} \in
\Hom_A(B', A) \otimes_{B'} B'_g = \omega_{B'_g/A} = \omega_{B/A}
$$
Nous omettons la vérification immédiate que cet élément ne dépend pas
du choix de $n$ et de $\psi$ ci-dessus.

\medskip\noindent
Montrons que $\tau$, construit au paragraphe précédent,
possède la propriété voulue dans un cas particulier. Supposons
$B' = C' \times D'$ et $g = (f, h)$, où $A \to C'$ est plat, $D'_h$ est plat, et
$f$ est inversible dans $C'$.
Il faut montrer que $\tau$ s'envoie sur $\text{Trace}_{C'/A}$ dans $\omega_{C'/A}$.
Dans ce cas, on choisit d'abord $n_D$ et
$\psi_D : D' \to \Hom_A(D', A) \otimes_A D'$ comme ci-dessus pour le couple
$(D', h)$, et l'on peut prendre pour
$\psi_C : C' \to \Hom_A(C', A) \otimes_A C' = \Hom_A(C', C')$
l'application envoyant $c' \in C'$ sur la multiplication par $c'$.
On prend alors $n = n_D$ et $\psi = (f^{n_D} \psi_C, \psi_D)$,
et la compatibilité voulue est claire car
$\text{Trace}_{C'/A} = ev \circ \psi_C$, comme on l'a remarqué ci-dessus.

\medskip\noindent
Pour établir la propriété voulue dans le cas général, supposons donnés
$A \to A_1$, avec $A_1$ noethérien, et une décomposition en produit
$B'_g \otimes_A A_1 = C \times D$ où $A_1 \to C$ est fini.
Posons $B'_1 = B' \otimes_A A_1$. Alors $\Spec(C) \to \Spec(B'_1)$
est une immersion ouverte puisque $B'_g \otimes_A A_1 = (B'_1)_g$, et
son image est fermée puisque $B'_1 \to C$ est fini (car $A_1 \to C$
est fini). Ainsi $B'_1 = C \times D'$ et $D'_g = D$.
On en déduit que $B'_1 = C \times D'$ et $g$, sur $A_1$,
sont dans la situation du paragraphe précédent.
Puisque la formation du diagramme affiché ci-dessus
commute au changement de base, la formation de $\tau$ commute
au changement de base $A \to A_1$ (détails omis ; utiliser la
résolution $F_1 \otimes_A A_1 \to F_0 \otimes_A A_1 \to B'_1 \to 0$
pour le voir). La compatibilité voulue découle donc du résultat
du paragraphe précédent.
\end{proof}

\begin{remark}
\label{remark-relative-dualizing-for-flat-quasi-finite}
Soit $f : Y \to X$ un morphisme plat localement quasi-fini de schémas localement
noethériens. Soit $\omega_{Y/X}$ comme dans la
remarque \ref{remark-relative-dualizing-for-quasi-finite}.
L'unicité, l'existence et la compatibilité à la
localisation des éléments trace
(lemmes \ref{lemma-trace-unique}, \ref{lemma-dualizing-tau} et
\ref{lemma-trace-base-change})
montrent qu'il existe une section globale
$$
\tau_{Y/X} \in \Gamma(Y, \omega_{Y/X})
$$
telle que, pour tout couple d'ouverts affines
$\Spec(B) = V \subset Y$, $\Spec(A) = U \subset X$ vérifiant $f(V) \subset U$,
cet élément $\tau_{Y/X}$ s'envoie sur $\tau_{B/A}$ par
l'isomorphisme canonique
$H^0(V, \omega_{Y/X}) = \omega_{B/A}$.
\end{remark}

\begin{lemma}
\label{lemma-tau-nonzero}
Soient $k$ un corps et $A$ une $k$-algèbre finie. Supposons $A$
locale, de corps résiduel $k'$. Les conditions suivantes sont équivalentes :
\begin{enumerate}
\item $\text{Trace}_{A/k}$ est non nulle ;
\item $\tau_{A/k} \in \omega_{A/k}$ est non nul ;
\item $k'/k$ est séparable et $\text{longueur}_A(A)$ est première
à la caractéristique de $k$.
\end{enumerate}
\end{lemma}

\begin{proof}
Les conditions (1) et (2) sont équivalentes d'après le lemme \ref{lemma-finite-flat-trace}.
Soit $\mathfrak m \subset A$. Puisque $\dim_k(A) < \infty$, il est clair que
$A$ est de longueur finie sur $A$. Choisissons une filtration
$$
A = I_0 \supset \mathfrak m = I_1 \supset I_2 \supset \ldots I_n = 0
$$
par des idéaux tels que $I_i/I_{i + 1} \cong k'$ comme $A$-modules. Voir
Algèbre, lemme \ref{algebra-lemma-simple-pieces}, qui montre également que
$n = \text{longueur}_A(A)$. Si $a \in \mathfrak m$, alors $aI_i \subset I_{i + 1}$,
et l'on voit immédiatement que $\text{Trace}_{A/k}(a) = 0$.
Si $a \not \in \mathfrak m$ a pour image $\lambda \in k'$, alors
on obtient
$$
\text{Trace}_{A/k}(a) =
\sum\nolimits_{i = 0, \ldots, n - 1}
\text{Trace}_k(a : I_i/I_{i - 1} \to I_i/I_{i - 1}) =
n \text{Trace}_{k'/k}(\lambda)
$$
On achève la démonstration en appliquant
Corps, lemme \ref{fields-lemma-separable-trace-pairing}.
\end{proof}






\section{Morphismes finis}
\label{section-finite-morphisms}

\noindent
Dans cette section, nous rassemblons quelques remarques sur les
constructions des sections précédentes dans le cas des morphismes finis.
Soit $f : Y \to X$ un morphisme fini de schémas localement noethériens.
Soit $\omega_{Y/X}$ comme dans la
remarque \ref{remark-relative-dualizing-for-quasi-finite}.

\medskip\noindent
La première remarque est que
$$
f_*\omega_{Y/X} = \SheafHom_{\mathcal{O}_X}(f_*\mathcal{O}_Y, \mathcal{O}_X)
$$
comme faisceaux de $f_*\mathcal{O}_Y$-modules. Puisque $f$ est affine, cette
formule caractérise $\omega_{Y/X}$ de façon unique, voir
Morphismes, lemme \ref{morphisms-lemma-affine-equivalence-modules}.
La formule est vraie car, pour un ouvert affine $\Spec(A) = U \subset X$,
l'image réciproque $V = f^{-1}(U)$ est le spectre d'une $A$-algèbre finie
$B$, d'où
$$
\begin{aligned}
H^0(U, f_*\omega_{Y/X})
&= H^0(V, \omega_{Y/X}) = \omega_{B/A} = \Hom_A(B, A) \\
&= H^0(U, \SheafHom_{\mathcal{O}_X}(f_*\mathcal{O}_Y, \mathcal{O}_X))
\end{aligned}
$$
par construction. On obtient en particulier une application canonique d'évaluation
$$
f_*\omega_{Y/X} \longrightarrow \mathcal{O}_X
$$
donnée par l'évaluation en $1$ lorsque l'on considère $f_*\omega_{Y/X}$
comme le faisceau $\SheafHom_{\mathcal{O}_X}(f_*\mathcal{O}_Y, \mathcal{O}_X)$.

\medskip\noindent
La deuxième remarque est que l'application d'évaluation fournit
des identifications canoniques
$$
\Hom_Y(\mathcal{F}, f^*\mathcal{G} \otimes_{\mathcal{O}_Y} \omega_{Y/X})
=
\Hom_X(f_*\mathcal{F}, \mathcal{G})
$$
fonctorielles en le module quasi-cohérent $\mathcal{F}$ sur $Y$
et en le module fini localement libre $\mathcal{G}$ sur $X$.
Si $\mathcal{G} = \mathcal{O}_X$, cela résulte immédiatement
de ce qui précède et de
Algèbre, lemme \ref{algebra-lemma-adjoint-hom-restrict}.
Pour $\mathcal{G}$ quelconque, on peut utiliser le même lemme et les
isomorphismes
$$
f_*(f^*\mathcal{G} \otimes_{\mathcal{O}_Y} \omega_{Y/X}) =
\mathcal{G} \otimes_{\mathcal{O}_X}
\SheafHom_{\mathcal{O}_X}(f_*\mathcal{O}_Y, \mathcal{O}_X) =
\SheafHom_{\mathcal{O}_X}(f_*\mathcal{O}_Y, \mathcal{G})
$$
de $f_*\mathcal{O}_Y$-modules, où la première égalité est la
formule de projection
(Cohomologie, lemme \ref{cohomology-lemma-projection-formula}).
On peut également démontrer la formule localement sur les ouverts affines par
un calcul direct.

\medskip\noindent
La troisième remarque est que, si $f$ est de plus plat, la
composée
$$
f_*\mathcal{O}_Y \xrightarrow{f_*\tau_{Y/X}} f_*\omega_{Y/X}
\longrightarrow \mathcal{O}_X
$$
est égale à l'application trace $\text{Trace}_f$ étudiée dans la
section \ref{section-discriminant}. Cela se voit immédiatement
sur les ouverts affines.

\medskip\noindent
La quatrième remarque est que, si $f$ est plat et $X$ noethérien, on
obtient
$$
\Hom_Y(K, Lf^*M \otimes_{\mathcal{O}_Y} \omega_{Y/X})
=
\Hom_X(Rf_*K, M)
$$
pour tous $K$ dans $D_\QCoh(\mathcal{O}_Y)$ et $M$ dans $D_\QCoh(\mathcal{O}_X)$.
Cela résulte des résultats de
Dualité des schémas, section \ref{duality-section-proper-flat},
mais on peut le démontrer directement dans ce cas, comme suit.
Tout d'abord, si $X$ est affine, l'assertion découle de
Complexes dualisants, lemmes \ref{dualizing-lemma-right-adjoint} et
\ref{dualizing-lemma-RHom-is-tensor-special}\footnote{Ce lemme admet
une démonstration plus simple dans notre cas.}, et de
Catégories dérivées de schémas, lemme \ref{perfect-lemma-affine-compare-bounded}.
On peut ensuite utiliser le principe de récurrence
(Cohomologie des schémas, lemme \ref{coherent-lemma-induction-principle})
et Mayer-Vietoris
(sous la forme de Cohomologie, lemme \ref{cohomology-lemma-mayer-vietoris-hom})
pour achever la démonstration.







\section{La différente de Noether}
\label{section-noether-different}

\noindent
La littérature contient de nombreuses notions de différente.
Nous en présentons quelques-unes dans cette section et les suivantes ; pour plus
de détails, nous renvoyons le lecteur à \cite{Kunz}.

\medskip\noindent
Soit $A \to B$ un homomorphisme d'anneaux. Notons
$$
\mu : B \otimes_A B \longrightarrow B,\quad
b \otimes b' \longmapsto bb'
$$
l'application de multiplication. Soit $I = \Ker(\mu)$. Il est clair que $I$ est
engendré par les éléments $b \otimes 1 - 1 \otimes b$ pour $b \in B$.
L'annulateur $J \subset B \otimes_A B$ de $I$ est donc canoniquement un $B$-module.
La {\it différente de Noether} de $B$ sur $A$ est
l'image de $J$ par l'application $\mu : B \otimes_A B \to B$. De façon équivalente,
la différente de Noether est l'image de l'application
$$
J = \Hom_{B \otimes_A B}(B, B \otimes_A B) \longrightarrow B,\quad
\varphi \longmapsto \mu(\varphi(1))
$$
Commençons par quelques lemmes indispensables.

\begin{lemma}
\label{lemma-noether-different-product}
Soient $A \to B_i$, $i = 1, 2$, des homomorphismes d'anneaux. Posons $B = B_1 \times B_2$.
\begin{enumerate}
\item L'annulateur $J$ de $\Ker(B \otimes_A B \to B)$ est $J_1 \times J_2$,
où $J_i$ est l'annulateur de $\Ker(B_i \otimes_A B_i \to B_i)$.
\item La différente de Noether $\mathfrak{D}$ de $B$ sur $A$ est
$\mathfrak{D}_1 \times \mathfrak{D}_2$, où $\mathfrak{D}_i$ est
la différente de Noether de $B_i$ sur $A$.
\end{enumerate}
\end{lemma}

\begin{proof}
Omise.
\end{proof}

\begin{lemma}
\label{lemma-noether-different-base-change}
Soit $A \to B$ un homomorphisme d'anneaux de type fini. Soit $A \to A'$ un homomorphisme plat d'anneaux.
Posons $B' = B \otimes_A A'$.
\begin{enumerate}
\item L'annulateur $J'$ de $\Ker(B' \otimes_{A'} B' \to B')$ est
$J \otimes_A A'$, où $J$ est l'annulateur de $\Ker(B \otimes_A B \to B)$.
\item La différente de Noether $\mathfrak{D}'$ de $B'$ sur $A'$ est
$\mathfrak{D}B'$, où $\mathfrak{D}$ est
la différente de Noether de $B$ sur $A$.
\end{enumerate}
\end{lemma}

\begin{proof}
Choisissons des générateurs $b_1, \ldots, b_n$ de $B$ comme $A$-algèbre.
Alors
$$
J = \Ker(B \otimes_A B \xrightarrow{b_i \otimes 1 - 1 \otimes b_i}
(B \otimes_A B)^{\oplus n})
$$
On voit donc que la formation de $J$ commute au changement de base plat.
Le résultat concernant la différente de Noether en découle immédiatement.
\end{proof}

\begin{lemma}
\label{lemma-noether-different-localization}
Soient $A \to B' \to B$ des homomorphismes d'anneaux, avec $A \to B'$
de type fini et $B' \to B$ induisant une immersion ouverte des spectres.
\begin{enumerate}
\item L'annulateur $J$ de $\Ker(B \otimes_A B \to B)$ est
$J' \otimes_{B'} B$, où $J'$ est l'annulateur de
$\Ker(B' \otimes_A B' \to B')$.
\item La différente de Noether $\mathfrak{D}$ de $B$ sur $A$ est
$\mathfrak{D}'B$, où $\mathfrak{D}'$ est
la différente de Noether de $B'$ sur $A$.
\end{enumerate}
\end{lemma}

\begin{proof}
Posons $I = \Ker(B \otimes_A B \to B)$ et $I' = \Ker(B' \otimes_A B' \to B')$.
Comme $\Spec(B) \to \Spec(B')$ est une immersion ouverte, on a
$B = (B \otimes_A B) \otimes_{B' \otimes_A B'} B'$. Il en résulte
$I = I'(B \otimes_A B)$. Puisque $I'$ est de type fini et que
$B' \otimes_A B' \to B \otimes_A B$ est plat, on obtient
$J = J'(B \otimes_A B)$, voir
Algèbre, lemme \ref{algebra-lemma-annihilator-flat-base-change}.
Puisque la structure de $B' \otimes_A B'$-module de $J'$
se factorise par $B' \otimes_A B' \to B'$, on en déduit (1).
L'assertion (2) est une conséquence de (1).
\end{proof}

\begin{remark}
\label{remark-construction-pairing}
Soit $A \to B$ un homomorphisme quasi-fini d'anneaux noethériens.
Soit $J$ l'annulateur de $\Ker(B \otimes_A B \to B)$.
Il existe un accouplement $B$-bilinéaire canonique
\begin{equation}
\label{equation-pairing-noether}
\omega_{B/A} \times J \longrightarrow B
\end{equation}
défini comme suit. Choisissons une factorisation $A \to B' \to B$
où $A \to B'$ est fini et $B' \to B$ induit une immersion ouverte
des spectres. Soit $J'$ l'annulateur de $\Ker(B' \otimes_A B' \to B')$.
On définit d'abord
$$
\Hom_A(B', A) \times J' \longrightarrow B',\quad
(\lambda, \sum b_i \otimes c_i) \longmapsto \sum \lambda(b_i)c_i
$$
Cet accouplement est $B'$-bilinéaire précisément parce que, pour $\xi \in J'$ et $b \in B'$,
on a $(b \otimes 1)\xi = (1 \otimes b)\xi$. D'après le
lemme \ref{lemma-noether-different-localization}
et puisque $\omega_{B/A} = \Hom_A(B', A) \otimes_{B'} B$,
on peut l'étendre en l'accouplement $B$-bilinéaire affiché ci-dessus.
\end{remark}

\begin{lemma}
\label{lemma-noether-pairing-compatibilities}
Soit $A \to B$ un homomorphisme quasi-fini d'anneaux noethériens.
\begin{enumerate}
\item Si $A \to A'$ est un homomorphisme plat d'anneaux noethériens, alors
$$
\xymatrix{
\omega_{B/A} \times J \ar[r] \ar[d] & B \ar[d]  \\
\omega_{B'/A'} \times J' \ar[r] & B'
}
$$
est commutatif, avec les notations du
lemme \ref{lemma-noether-different-base-change},
les flèches horizontales étant données par
(\ref{equation-pairing-noether}).
\item Si $B = B_1 \times B_2$, alors
$$
\xymatrix{
\omega_{B/A} \times J \ar[r] \ar[d] & B \ar[d]  \\
\omega_{B_i/A} \times J_i \ar[r] & B_i
}
$$
est commutatif pour $i = 1, 2$, avec les notations du
lemme \ref{lemma-noether-different-product},
les flèches horizontales étant données par
(\ref{equation-pairing-noether}).
\end{enumerate}
\end{lemma}

\begin{proof}
En vertu de la construction de l'accouplement dans la
remarque \ref{remark-construction-pairing},
(1) et (2) se ramènent au cas où $A \to B$ est fini.
L'assertion (1) résulte alors de ce que l'application de contraction
$\Hom_A(M, A) \otimes_A M \otimes_A M \to M$,
$\lambda \otimes m \otimes m' \mapsto \lambda(m)m'$,
commute au changement de base. Pour (2), utiliser le fait que
$J = J_1 \times J_2$ est contenu dans les facteurs
$B_1 \otimes_A B_1$ et $B_2 \otimes_A B_2$
de $B \otimes_A B$.
\end{proof}

\begin{lemma}
\label{lemma-noether-pairing-flat-quasi-finite}
{\emergencystretch=2em Soit $A \to B$ un homomorphisme plat quasi-fini d'anneaux noethériens.
L'accouplement de la remarque \ref{remark-construction-pairing} induit un isomorphisme
$J \to \Hom_B(\omega_{B/A}, B)$.\par}
\end{lemma}

\begin{proof}
Montrons d'abord l'assertion lorsque $A \to B$ est fini plat. Dans ce cas, on peut
localiser sur $A$ et supposer que $B$ est un $A$-module libre de type fini. Soit
$b_1, \ldots, b_n$ une base de $B$ comme $A$-module, et notons
$b_1^\vee, \ldots, b_n^\vee$ la base duale de $\omega_{B/A}$. Remarquons que
$\sum b_i \otimes c_i \in J$ s'envoie sur l'élément de $\Hom_B(\omega_{B/A}, B)$
qui envoie $b_i^\vee$ sur $c_i$. Supposons $\varphi : \omega_{B/A} \to B$
$B$-linéaire. Nous affirmons alors que $\xi = \sum b_i \otimes \varphi(b_i^\vee)$
est un élément de $J$. En effet, la $B$-linéarité de $\varphi$
entraîne précisément que $(b \otimes 1)\xi = (1 \otimes b)\xi$ pour tout $b \in B$.
Notre application possède donc un inverse et est un isomorphisme.

\medskip\noindent
Soit $\mathfrak q \subset B$ un idéal premier au-dessus de $\mathfrak p \subset A$.
Nous allons montrer que l'application localisée
$$
J_\mathfrak q
\longrightarrow
\Hom_B(\omega_B/A, B)_\mathfrak q
$$
est un isomorphisme.
Cela suffit d'après Algèbre, lemme \ref{algebra-lemma-characterize-zero-local}.
D'après
Algèbre, lemme \ref{algebra-lemma-etale-makes-quasi-finite-finite-one-prime},
on peut trouver un homomorphisme étale d'anneaux $A \to A'$ et un idéal
premier $\mathfrak p' \subset A'$ au-dessus de $\mathfrak p$
tels que $\kappa(\mathfrak p') = \kappa(\mathfrak p)$ et
que
$$
B' = B \otimes_A A' = C \times D
$$
avec $A' \to C$ fini, et tels que l'unique idéal premier $\mathfrak q'$
de $B \otimes_A A'$ au-dessus de $\mathfrak q$ et de $\mathfrak p'$
corresponde à un idéal premier de $C$. Soit $J'$ l'annulateur de
$\Ker(B' \otimes_{A'} B' \to B')$. D'après les
lemmes \ref{lemma-dualizing-flat-base-change},
\ref{lemma-noether-different-base-change} et
\ref{lemma-noether-pairing-compatibilities},
l'application $J' \to \Hom_{B'}(\omega_{B'/A'}, B')$
s'obtient en appliquant le foncteur $- \otimes_B B'$
à l'application $J \to \Hom_B(\omega_{B/A}, B)$.
Puisque $B_\mathfrak q \to B'_{\mathfrak q'}$ est fidèlement plat,
il suffit d'établir le résultat pour $(A' \to B', \mathfrak q')$.
D'après les lemmes \ref{lemma-dualizing-product},
\ref{lemma-noether-different-product} et
\ref{lemma-noether-pairing-compatibilities},
on est ramené au cas traité dans le premier
paragraphe de la démonstration.
\end{proof}

\begin{lemma}
\label{lemma-noether-different-flat-quasi-finite}
Soit $A \to B$ un homomorphisme plat quasi-fini d'anneaux noethériens.
Le diagramme
$$
\xymatrix{
J \ar[rr] \ar[rd]_\mu & &
\Hom_B(\omega_{B/A}, B) \ar[ld]^{\varphi \mapsto \varphi(\tau_{B/A})} \\
& B
}
$$
commute, la flèche horizontale étant l'isomorphisme du
lemme \ref{lemma-noether-pairing-flat-quasi-finite}.
La différente de Noether de $B$ sur $A$
est donc l'image de l'application $\Hom_B(\omega_{B/A}, B) \to B$.
\end{lemma}

\begin{proof}
Exactement comme dans la démonstration du lemme \ref{lemma-noether-pairing-flat-quasi-finite},
on se ramène au cas d'un homomorphisme fini libre $A \to B$.
Dans ce cas, $\tau_{B/A} = \text{Trace}_{B/A}$.
Choisissons une base $b_1, \ldots, b_n$ de $B$ comme $A$-module.
Soit $\xi = \sum b_i \otimes c_i \in J$. Alors $\mu(\xi) = \sum b_i c_i$.
D'autre part, l'image de $\xi$ dans $\Hom_B(\omega_{B/A}, B)$
envoie $\text{Trace}_{B/A}$ sur $\sum \text{Trace}_{B/A}(b_i)c_i$.
Il faut donc montrer que
$$
\sum b_ic_i = \sum \text{Trace}_{B/A}(b_i)c_i
$$
lorsque $\xi = \sum b_i \otimes c_i \in J$. Écrivons $b_i b_j = \sum_k a_{ij}^k b_k$
avec $a_{ij}^k \in A$. Le membre de droite est alors
$\sum_{i, j} a_{ij}^j c_i$. D'autre part, $\xi \in J$ entraîne
$$
(b_j \otimes 1)(\sum\nolimits_i b_i \otimes c_i) =
(1 \otimes b_j)(\sum\nolimits_i b_i \otimes c_i)
$$
ce qui donne $b_j c_i = \sum_k a_{jk}^i c_k$. Le membre de gauche
est donc $\sum_{i, j} a_{ij}^i c_j$. Puisque $a_{ij}^k = a_{ji}^k$, on a bien l'égalité.
\end{proof}

\begin{lemma}
\label{lemma-noether-different}
Soit $A \to B$ un homomorphisme d'anneaux de type fini. Soit $\mathfrak{D} \subset B$
la différente de Noether. Alors $V(\mathfrak{D})$ est l'ensemble des idéaux premiers
$\mathfrak q \subset B$ tels que $A \to B$ soit ramifié en $\mathfrak q$.
\end{lemma}

\begin{proof}
Supposons $A \to B$ non ramifié en $\mathfrak q$. En remplaçant
$B$ par $B_g$ pour un $g \in B$, $g \not \in \mathfrak q$, on peut
supposer $A \to B$ non ramifié (Algèbre, définition
\ref{algebra-definition-unramified} et
lemme \ref{lemma-noether-different-localization}).
Dans ce cas, $\Omega_{B/A} = 0$. Ainsi, si $I = \Ker(B \otimes_A B \to B)$,
alors $I/I^2 = 0$ d'après
Algèbre, lemme \ref{algebra-lemma-differentials-diagonal}.
Puisque $A \to B$ est de type fini, $I$ est de type
fini. Le lemme de Nakayama
(Algèbre, lemme \ref{algebra-lemma-NAK})
assure donc l'existence d'un élément de la forme $1 + i$
annulant $I$. Il en résulte que $\mathfrak{D} = B$.

\medskip\noindent
Réciproquement, supposons $\mathfrak{D} \not \subset \mathfrak q$.
En remplaçant $B$ par une localisation principale comme ci-dessus,
on peut supposer $\mathfrak{D} = B$. Cela signifie qu'il existe un
élément de la forme $1 + i$ dans l'annulateur de $I$.
Réciproquement, cela entraîne que $I/I^2 = \Omega_{B/A}$ est nul,
ce qui permet de conclure.
\end{proof}







\section{La différente de Kähler}
\label{section-kahler-different}

\noindent
Soit $A \to B$ un homomorphisme d'anneaux de type fini. La {\it différente de Kähler} est
l'idéal de Fitting d'ordre zéro de $\Omega_{B/A}$ comme $B$-module. On globalise cette
définition comme suit.

\begin{definition}
\label{definition-kahler-different}
Soit $f : Y \to X$ un morphisme de schémas localement de type fini.
La {\it différente de Kähler} est l'idéal de Fitting d'ordre $0$ de $\Omega_{Y/X}$.
\end{definition}

\noindent
La différente de Kähler est un faisceau quasi-cohérent d'idéaux sur $Y$.

\begin{lemma}
\label{lemma-base-change-kahler-different}
Considérons un diagramme cartésien de schémas
$$
\xymatrix{
Y' \ar[d]_{f'} \ar[r] & Y \ar[d]^f \\
X' \ar[r]^g & X
}
$$
où $f$ est localement de type fini. Soit $R \subset Y$, resp.\ $R' \subset Y'$,
le sous-schéma fermé défini par la différente de Kähler de $f$, resp.\ de $f'$.
Alors $Y' \to Y$ induit un isomorphisme $R' \to R \times_Y Y'$.
\end{lemma}

\begin{proof}
Cela tient au fait que $\Omega_{Y'/X'}$ est l'image réciproque de $\Omega_{Y/X}$
(Morphismes, lemme \ref{morphisms-lemma-base-change-differentials}) ;
on peut alors appliquer
Compléments d'algèbre, lemme \ref{more-algebra-lemma-fitting-ideal-basics}.
\end{proof}

\begin{lemma}
\label{lemma-kahler-different}
Soit $f : Y \to X$ un morphisme de schémas localement de type fini.
Soit $R \subset Y$ le sous-schéma fermé défini par
la différente de Kähler. Alors $R \subset Y$ est exactement
l'ensemble des points où $f$ est ramifié.
\end{lemma}

\begin{proof}
Il s'agit de l'énoncé de
Diviseurs, lemme \ref{divisors-lemma-zero-fitting-ideal-omega-unramified}.
\end{proof}

\begin{lemma}
\label{lemma-kahler-different-complete-intersection}
Soient $A$ un anneau, $n \geq 1$ et
$f_1, \ldots, f_n \in A[x_1, \ldots, x_n]$.
Posons $B = A[x_1, \ldots, x_n]/(f_1, \ldots, f_n)$.
La différente de Kähler de $B$ sur $A$ est l'idéal
de $B$ engendré par $\det(\partial f_i/\partial x_j)$.
\end{lemma}

\begin{proof}
En effet, $\Omega_{B/A}$ admet la présentation
$$
\bigoplus\nolimits_{i = 1, \ldots, n} B f_i
\xrightarrow{\text{d}}
\bigoplus\nolimits_{j = 1, \ldots, n} B \text{d}x_j
\rightarrow \Omega_{B/A} \rightarrow 0
$$
d'après Algèbre, lemme \ref{algebra-lemma-differential-seq}.
\end{proof}



\section{La différente de Dedekind}
\label{section-dedekind-different}

\noindent
Soit $A \to B$ un homomorphisme d'anneaux. On dit que {\it la différente de Dedekind est définie}
si $A$ est noethérien, $A \to B$ est fini,
tout non-diviseur de zéro sur $A$ est non diviseur de zéro sur $B$, et $K \to L$ est
étale, où $K = Q(A)$ et $L = B \otimes_A K$. Alors $K \subset L$ est
fini étale et
$$
\mathcal{L}_{B/A} = \{x \in L \mid \text{Trace}_{L/K}(bx) \in A
\text{ pour tout }b \in B\}
$$
est le module complémentaire de Dedekind. Dans cette situation, la
{\it différente de Dedekind} est
$$
\mathfrak{D}_{B/A} = \{x \in L \mid x\mathcal{L}_{B/A} \subset B\}
$$
considérée comme un $B$-sous-module de $L$.
D'après le lemme \ref{lemma-dedekind-different-ideal}, la différente de Dedekind est un
idéal de $B$ si $A$ est normal ou si $B$ est plat sur $A$.

\begin{lemma}
\label{lemma-dedekind-different-ideal}
Supposons que la différente de Dedekind de $A \to B$ soit définie. Considérons les assertions
\begin{enumerate}
\item $A \to B$ est plat ;
\item $A$ est un anneau normal ;
\item $\text{Trace}_{L/K}(B) \subset A$ ;
\item $1 \in \mathcal{L}_{B/A}$ ;
\item la différente de Dedekind $\mathfrak{D}_{B/A}$ est un idéal de $B$.
\end{enumerate}
On a alors (1) $\Rightarrow$ (3), (2) $\Rightarrow$ (3),
(3) $\Leftrightarrow$ (4) et (4) $\Rightarrow$ (5).
\end{lemma}

\begin{proof}
L'équivalence de (3) et (4) ainsi que
l'implication (4) $\Rightarrow$ (5) sont immédiates.

\medskip\noindent
Si $A \to B$ est plat, $\text{Trace}_{B/A} : B \to A$ est
définie et $\text{Trace}_{L/K}$ s'en déduit par changement de base. On a donc (3).

\medskip\noindent
Si $A$ est normal, $A$ est un produit fini d'anneaux intègres normaux,
et l'on se ramène donc au cas d'un anneau intègre normal. Alors $K$ est
le corps des fractions de $A$, et $L = \prod L_i$ est un produit fini
d'extensions finies séparables de $K$. On a
$\text{Trace}_{L/K}(b) = \sum \text{Trace}_{L_i/K}(b_i)$,
où $b_i \in L_i$ est l'image de $b$.
Puisque $b$ est entier sur $A$, car $B$ est fini sur $A$,
ces traces appartiennent à $A$. En effet, le
polynôme minimal de $b_i$ sur $K$ a ses coefficients dans $A$
(Algèbre, lemme \ref{algebra-lemma-minimal-polynomial-normal-domain}),
et $\text{Trace}_{L_i/K}(b_i)$ est un
multiple entier de l'un de ces coefficients
(Corps, lemme \ref{fields-lemma-trace-and-norm-from-minimal-polynomial}).
\end{proof}

\begin{lemma}
\label{lemma-dedekind-complementary-module}
Si la différente de Dedekind de $A \to B$ est définie, alors
il existe un isomorphisme canonique
$\mathcal{L}_{B/A} \to \omega_{B/A}$.
\end{lemma}

\begin{proof}
Rappelons que $\omega_{B/A} = \Hom_A(B, A)$ puisque $A \to B$ est fini.
On envoie $x \in \mathcal{L}_{B/A}$ sur l'application
$b \mapsto \text{Trace}_{L/K}(bx)$.
Réciproquement, une application $A$-linéaire $\varphi : B \to A$
fournit une application $K$-linéaire $\varphi_K : L \to K$. Puisque $K \to L$ est fini
étale, la forme trace est non dégénérée
(lemme \ref{lemma-discriminant}), et il existe donc $x \in L$ tel que
$\varphi_K(y) = \text{Trace}_{L/K}(xy)$ pour tout $y \in L$.
Alors $x \in \mathcal{L}_{B/A}$ s'envoie sur $\varphi$ dans $\omega_{B/A}$.
\end{proof}

\begin{lemma}
\label{lemma-flat-dedekind-complementary-module-trace}
Si la différente de Dedekind de $A \to B$ est définie et si $A \to B$ est plat, alors
\begin{enumerate}
\item l'isomorphisme canonique $\mathcal{L}_{B/A} \to \omega_{B/A}$
envoie $1 \in \mathcal{L}_{B/A}$ sur l'élément trace
$\tau_{B/A} \in \omega_{B/A}$ ;
\item la différente de Dedekind est
$\mathfrak{D}_{B/A} = \{b \in B \mid b\omega_{B/A} \subset B\tau_{B/A}\}$.
\end{enumerate}
\end{lemma}

\begin{proof}
La première assertion
résulte de la démonstration du lemme \ref{lemma-dedekind-different-ideal}
et du lemme \ref{lemma-finite-flat-trace}.
La deuxième découle immédiatement de la première et des
définitions.
\end{proof}



\section{La différente}
\label{section-different}

\noindent
La définition suivante est motivée par le fait qu'elle redonne la
différente de Dedekind dans le cas fini plat, comme nous le verrons ci-dessous.

\begin{definition}
\label{definition-different}
Soit $f : Y \to X$ un morphisme plat localement quasi-fini de
schémas localement noethériens.
Soient $\omega_{Y/X}$ le module dualisant relatif et
$\tau_{Y/X} \in \Gamma(Y, \omega_{Y/X})$ l'élément trace
(remarques \ref{remark-relative-dualizing-for-quasi-finite} et
\ref{remark-relative-dualizing-for-flat-quasi-finite}).
L'annulateur de
$$
\Coker(\mathcal{O}_Y \xrightarrow{\tau_{Y/X}} \omega_{Y/X})
$$
est la {\it différente} de $Y/X$. C'est un idéal cohérent
$\mathfrak{D}_f \subset \mathcal{O}_Y$.
\end{definition}

\noindent
Nous généraliserons ceci dans la remarque \ref{remark-different-generalization} ci-dessous.
Remarquons que $\mathfrak{D}_f$ est localement engendré par un élément si
$\omega_{Y/X}$ est un $\mathcal{O}_Y$-module inversible.
Énonçons d'abord l'égalité avec la différente de Dedekind.

\begin{lemma}
\label{lemma-flat-agree-dedekind}
Soit $f : Y \to X$ un morphisme plat quasi-fini de schémas noethériens.
Soient $V = \Spec(B) \subset Y$, $U = \Spec(A) \subset X$
des sous-schémas ouverts affines tels que $f(V) \subset U$.
Si la différente de Dedekind de $A \to B$ est définie, alors
$$
\mathfrak{D}_f|_V = \widetilde{\mathfrak{D}_{B/A}}
$$
comme faisceaux cohérents d'idéaux sur $V$.
\end{lemma}

\begin{proof}
Cela résulte immédiatement des lemmes \ref{lemma-dedekind-different-ideal} et
\ref{lemma-flat-dedekind-complementary-module-trace}.
\end{proof}

\begin{lemma}
\label{lemma-flat-gorenstein-agree-noether}
Soit $f : Y \to X$ un morphisme plat quasi-fini de schémas noethériens.
Soient $V = \Spec(B) \subset Y$, $U = \Spec(A) \subset X$
des sous-schémas ouverts affines tels que $f(V) \subset U$.
Si $\omega_{Y/X}|_V$ est inversible, c'est-à-dire si $\omega_{B/A}$
est un $B$-module inversible, alors
$$
\mathfrak{D}_f|_V = \widetilde{\mathfrak{D}}
$$
comme faisceaux cohérents d'idéaux sur $V$, où
$\mathfrak{D} \subset B$ est la différente de Noether de $B$ sur $A$.
\end{lemma}

\begin{proof}
Considérons l'application
$$
\SheafHom_{\mathcal{O}_Y}(\omega_{Y/X}, \mathcal{O}_Y)
\longrightarrow
\mathcal{O}_Y,\quad
\varphi \longmapsto \varphi(\tau_{Y/X})
$$
Son image correspond à la différente de Noether
sur les ouverts affines, voir le lemme \ref{lemma-noether-different-flat-quasi-finite}.
Le résultat découle donc du fait élémentaire suivant : étant donné
un module inversible $\omega$ et une section globale $\tau$,
l'image de
$\tau : \SheafHom(\omega, \mathcal{O}) = \omega^{\otimes -1} \to \mathcal{O}$
est égale à l'annulateur de $\Coker(\tau : \mathcal{O} \to \omega)$.
\end{proof}

\begin{lemma}
\label{lemma-base-change-different}
Considérons un diagramme cartésien de schémas noethériens
$$
\xymatrix{
Y' \ar[d]_{f'} \ar[r] & Y \ar[d]^f \\
X' \ar[r]^g & X
}
$$
où $f$ est plat et quasi-fini. Soit $R \subset Y$, resp.\ $R' \subset Y'$,
le sous-schéma fermé défini par la différente
$\mathfrak{D}_f$, resp.\ $\mathfrak{D}_{f'}$.
Alors $Y' \to Y$ induit une immersion fermée bijective $R' \to R \times_Y Y'$.
Si $g$ est plat ou si $\omega_{Y/X}$ est inversible, alors
$R' = R \times_Y Y'$.
\end{lemma}

\begin{proof}
On se ramène immédiatement au cas où $X$, $X'$, $Y$, $Y'$
sont affines. Autrement dit, on dispose d'un diagramme cocartésien d'anneaux
noethériens
$$
\xymatrix{
B' & B \ar[l] \\
A' \ar[u] & A \ar[l] \ar[u]
}
$$
où $A \to B$ est plat et quasi-fini. L'application de changement de base
$\omega_{B/A} \otimes_B B' \to \omega_{B'/A'}$ est un isomorphisme
(lemme \ref{lemma-dualizing-base-change-of-flat}) et envoie
l'élément trace $\tau_{B/A}$ sur l'élément trace $\tau_{B'/A'}$
(lemme \ref{lemma-trace-base-change}).
Le $B$-module de type fini $Q = \Coker(\tau_{B/A} : B \to \omega_{B/A})$
vérifie donc $Q \otimes_B B' = \Coker(\tau_{B'/A'} : B' \to \omega_{B'/A'})$.
Ainsi $\mathfrak{D}_{B/A}B' \subset \mathfrak{D}_{B'/A'}$, ce qui signifie
que l'on obtient l'immersion fermée $R' \to R \times_Y Y'$.
Puisque $R = \text{Supp}(Q)$ et $R' = \text{Supp}(Q \otimes_B B')$
(Algèbre, lemme \ref{algebra-lemma-support-closed}),
on voit que $R' \to R \times_Y Y'$ est bijective d'après
Algèbre, lemme \ref{algebra-lemma-support-base-change}.
L'égalité $\mathfrak{D}_{B/A}B' = \mathfrak{D}_{B'/A'}$ est vraie
si $B \to B'$ est plat, par exemple si $A \to A'$ est plat, voir
Algèbre, lemme \ref{algebra-lemma-annihilator-flat-base-change}.
Enfin, si $\omega_{B/A}$ est inversible, on peut localiser
et supposer $\omega_{B/A} = B \lambda$. En écrivant $\tau_{B/A} = b\lambda$,
on voit que $Q = B/bB$ et $\mathfrak{D}_{B/A} = bB$.
Le même raisonnement sur $B'$
donne $\mathfrak{D}_{B'/A'} = bB'$, ce qui démontre le lemme.
\end{proof}

\begin{lemma}
\label{lemma-norm-different-in-discriminant}
Soit $f : Y \to X$ un morphisme fini plat de schémas noethériens.
Alors $\text{Norm}_f : f_*\mathcal{O}_Y \to \mathcal{O}_X$ envoie
$f_*\mathfrak{D}_f$ dans le faisceau d'idéaux du discriminant $D_f$.
\end{lemma}

\begin{proof}
L'application norme est construite dans
Diviseurs, lemme \ref{divisors-lemma-finite-locally-free-has-norm},
et le discriminant de $f$ dans la section \ref{section-discriminant}.
La question est locale sur les ouverts affines ; on peut donc supposer $X = \Spec(A)$,
$Y = \Spec(B)$ et $f$ donné par un homomorphisme fini localement libre d'anneaux $A \to B$.
En localisant davantage, on peut supposer $B$ libre de type fini comme $A$-module.
Choisissons une base $b_1, \ldots, b_n \in B$ de $B$ comme $A$-module.
Notons $b_1^\vee, \ldots, b_n^\vee$ la base duale de
$\omega_{B/A} = \Hom_A(B, A)$ comme $A$-module.
Puisque la norme de $b$ est le déterminant de $b : B \to B$ comme
application $A$-linéaire, on a
$\text{Norm}_{B/A}(b) = \det(b_i^\vee(bb_j))$.
Le discriminant est le sous-schéma fermé principal de $\Spec(A)$
défini par $\det(\text{Trace}_{B/A}(b_ib_j))$.
Si $b \in \mathfrak{D}_{B/A}$, alors
il existe $c_i \in B$ tels que
$b \cdot b_i^\vee = c_i \cdot \text{Trace}_{B/A}$, où
le point désigne la structure de $B$-module sur $\omega_{B/A}$.
Écrivons $c_i = \sum a_{il} b_l$.
On a
\begin{align*}
\text{Norm}_{B/A}(b)
& =
\det(b_i^\vee(bb_j)) \\
& =
\det( (b \cdot b_i^\vee)(b_j)) \\
& =
\det((c_i \cdot \text{Trace}_{B/A})(b_j)) \\
& =
\det(\text{Trace}_{B/A}(c_ib_j)) \\
& =
\det(a_{il}) \det(\text{Trace}_{B/A}(b_l b_j))
\end{align*}
ce qui démontre le lemme.
\end{proof}

\begin{lemma}
\label{lemma-different-ramification}
Soit $f : Y \to X$ un morphisme plat quasi-fini de schémas noethériens.
Le sous-schéma fermé $R \subset Y$ défini par la différente $\mathfrak{D}_f$
est exactement l'ensemble des points où $f$ n'est pas étale
(ou, de façon équivalente, est ramifié).
\end{lemma}

\begin{proof}
Puisque $f$ est de présentation finie et plat, il est étale
en un point si et seulement s'il est non ramifié en ce point. De plus, la
formation du lieu des points ramifiés commute au changement de base.
Voir Morphismes, section \ref{morphisms-section-etale}, et notamment
Morphismes, lemme \ref{morphisms-lemma-set-points-where-fibres-etale}.
D'après le lemme \ref{lemma-base-change-different}, la formation de $R$ commute,
ensemblistement, au changement de base. Il suffit donc de démontrer le
lemme lorsque $X$ est le spectre d'un corps. D'autre part, la
construction de $(\omega_{Y/X}, \tau_{Y/X})$ est locale sur $Y$.
Puisque $Y$ est un espace discret fini (étant quasi-fini
sur un corps), on peut supposer que $Y$ possède un unique point.

\medskip\noindent
Écrivons $X = \Spec(k)$ et $Y = \Spec(B)$, où $k$ est un corps et $B$ une
$k$-algèbre locale finie. Si $Y \to X$ est étale, alors
$B$ est une extension finie séparable de $k$, et l'élément
trace $\text{Trace}_{B/k}$ est un élément de base de $\omega_{B/k}$
d'après Corps, lemme \ref{fields-lemma-separable-trace-pairing}.
Ainsi $\mathfrak{D}_{B/k} = B$ dans ce cas.
Réciproquement, si $\mathfrak{D}_{B/k} = B$, le
lemme \ref{lemma-norm-different-in-discriminant}
et le fait que la norme de $1$ est $1$ montrent que le
discriminant est vide. Par conséquent,
$Y \to X$ est étale d'après le lemme \ref{lemma-discriminant}.
\end{proof}

\begin{lemma}
\label{lemma-norm-different-is-discriminant}
Soit $f : Y \to X$ un morphisme plat quasi-fini de schémas noethériens.
Soit $R \subset Y$ le sous-schéma fermé défini par $\mathfrak{D}_f$.
\begin{enumerate}
\item Si $\omega_{Y/X}$ est inversible,
alors $R$ est un sous-schéma fermé localement principal de $Y$.
\item Si $\omega_{Y/X}$ est inversible et $f$ fini, alors
la norme de $R$ est le discriminant $D_f$ de $f$.
\item Si $\omega_{Y/X}$ est inversible et $f$
est étale aux points associés de $Y$, alors $R$
est un diviseur de Cartier effectif et l'on a un
isomorphisme $\mathcal{O}_Y(R) = \omega_{Y/X}$.
\end{enumerate}
\end{lemma}

\begin{proof}
Preuve de (1). On peut travailler localement sur $Y$, donc supposer
$\omega_{Y/X}$ libre de rang $1$. Écrivons $\omega_{Y/X} = \mathcal{O}_Y\lambda$.
On peut alors écrire $\tau_{Y/X} = h \lambda$, et l'on voit que
$R$ est défini par $h$, c'est-à-dire que $R$ est localement principal.

\medskip\noindent
Preuve de (2). On peut supposer que $Y \to X$ est donné par un homomorphisme fini libre
d'anneaux $A \to B$ et que $\omega_{B/A}$ est libre de rang $1$ comme $B$-module.
Choisissons une $B$-base $\lambda$ de $\omega_{B/A}$ et écrivons
$\text{Trace}_{B/A} = b \cdot \lambda$ avec $b \in B$.
Alors $\mathfrak{D}_{B/A} = (b)$, et $D_f$ est défini par
$\det(\text{Trace}_{B/A}(b_ib_j))$, où $b_1, \ldots, b_n$ est une
base de $B$ comme $A$-module. Soit $b_1^\vee, \ldots, b_n^\vee$
la base duale.
En écrivant $b_i^\vee = c_i \cdot \lambda$, on voit que
$c_1, \ldots, c_n$ est également une base de $B$.
Ainsi, en posant $c_i = \sum a_{il}b_l$, on voit que $\det(a_{il})$
est inversible dans $A$. Il est clair que
$b \cdot b_i^\vee = c_i \cdot \text{Trace}_{B/A}$ ;
on conclut donc du calcul fait dans la démonstration du
lemme \ref{lemma-norm-different-in-discriminant}
que $\text{Norm}_{B/A}(b)$ est le produit d'une unité par
$\det(\text{Trace}_{B/A}(b_ib_j))$.

\medskip\noindent
Preuve de (3). Avec les notations ci-dessus, le
lemme \ref{lemma-different-ramification} et l'hypothèse montrent
que $h$ ne s'annule pas aux
points associés de $Y$, ce qui entraîne que $h$ est non diviseur de zéro.
L'isomorphisme canonique envoie $1$ sur $\tau_{Y/X}$, voir
Diviseurs, lemme \ref{divisors-lemma-characterize-OD}.
\end{proof}







\section{Morphismes quasi-finis syntomiques}
\label{section-quasi-finite-syntomic}

\noindent
Cette section établit qu'un morphisme quasi-fini syntomique
possède un module dualisant relatif inversible.

\begin{lemma}
\label{lemma-syntomic-quasi-finite}
Soit $f : Y \to X$ un morphisme de schémas. Les conditions suivantes sont équivalentes :
\begin{enumerate}
\item $f$ est localement quasi-fini et syntomique ;
\item $f$ est localement quasi-fini, plat et localement d'intersection
complète ;
\item $f$ est localement quasi-fini, plat, localement de présentation finie,
et les fibres de $f$ sont localement des intersections complètes ;
\item $f$ est localement quasi-fini et, pour tout $y \in Y$, il existe des
ouverts affines $y \in V = \Spec(B) \subset Y$, $U = \Spec(A) \subset X$
avec $f(V) \subset U$, un entier $n$ et
$h, f_1, \ldots, f_n \in A[x_1, \ldots, x_n]$ tels que
$B = A[x_1, \ldots, x_n, 1/h]/(f_1, \ldots, f_n)$ ;
\item pour tout $y \in Y$, il existe des ouverts affines
$y \in V = \Spec(B) \subset Y$, $U = \Spec(A) \subset X$
avec $f(V) \subset U$ tels que $A \to B$ soit une intersection complète globale
relative de la forme $B = A[x_1, \ldots, x_n]/(f_1, \ldots, f_n)$ ;
\item $f$ est localement quasi-fini, plat, localement de présentation finie,
et $\NL_{Y/X}$ est de Tor-amplitude contenue dans $[-1, 0]$ ;
\item $f$ est plat, localement de présentation finie,
$\NL_{Y/X}$ est parfait de rang $0$ et de Tor-amplitude contenue dans $[-1, 0]$.
\end{enumerate}
\end{lemma}

\begin{proof}
L'équivalence de (1) et (2) est donnée par
Compléments sur les morphismes, lemme \ref{more-morphisms-lemma-flat-lci}.
L'équivalence de (1) et (3) est donnée par
Morphismes, lemme \ref{morphisms-lemma-syntomic-flat-fibres}.

\medskip\noindent
Si $A \to B$ est comme dans (4), alors
$B = A[x, x_1, \ldots, x_n]/(xh - 1, f_1, \ldots, f_n)$
est une intersection complète globale relative, voir Algèbre, définition
\ref{algebra-definition-relative-global-complete-intersection}.
Ainsi (4) entraîne (5).
Il est clair que (5) entraîne (4).

\medskip\noindent
La condition (5) entraîne (1) : d'après
Algèbre, lemme \ref{algebra-lemma-relative-global-complete-intersection},
une intersection complète globale relative est syntomique, et
la définition d'une intersection complète globale relative
garantit qu'une telle intersection à
$n$ variables et $n$ équations est quasi-finie, voir
Algèbre, définition
\ref{algebra-definition-relative-global-complete-intersection} et
lemme \ref{algebra-lemma-isolated-point-fibre}.

\medskip\noindent
Algèbre, lemme \ref{algebra-lemma-syntomic}, ou
Morphismes, lemme \ref{morphisms-lemma-syntomic-locally-standard-syntomic},
montre que (1) entraîne (5).

\medskip\noindent
Compléments sur les morphismes, lemme \ref{more-morphisms-lemma-flat-fp-NL-lci}, montre que
(6) équivaut à (1). Si les conditions équivalentes (1) -- (6) sont satisfaites,
on voit que, localement sur les ouverts affines, $Y \to X$ est donné par une intersection complète
globale relative $B = A[x_1, \ldots, x_n]/(f_1, \ldots, f_n)$
ayant autant de variables que
d'équations. Cette présentation donne
$$
\NL_{B/A} =\left(
(f_1, \ldots, f_n)/(f_1, \ldots, f_n)^2
\longrightarrow
\bigoplus\nolimits_{i = 1, \ldots, n} B \text{d} x_i\right)
$$
D'après Algèbre, lemme
\ref{algebra-lemma-relative-global-complete-intersection-conormal},
le module $(f_1, \ldots, f_n)/(f_1, \ldots, f_n)^2$
est libre, de base formée des classes des éléments
$f_1, \ldots, f_n$. Ainsi $\NL_{B/A}$ est de rang $0$, de même que $\NL_{Y/X}$.
On voit ainsi que (1) -- (6) entraînent (7).

\medskip\noindent
Enfin, supposons (7). D'après
Compléments sur les morphismes, lemme \ref{more-morphisms-lemma-flat-fp-NL-lci},
le morphisme $f$ est syntomique. Sur des ouverts affines convenables,
$f$ est donc donné par une intersection complète globale relative
$A \to B = A[x_1, \ldots, x_n]/(f_1, \ldots, f_m)$, voir
Morphismes, lemme \ref{morphisms-lemma-syntomic-locally-standard-syntomic}.
Exactement comme ci-dessus, on voit que $\NL_{B/A}$ est un complexe parfait
de rang $n - m$. Ainsi $n = m$ et la condition (5) est satisfaite.
La démonstration est achevée.
\end{proof}

\begin{lemma}
\label{lemma-characterize-invertible}
Inversibilité du module dualisant relatif.
\begin{enumerate}
\item Si $A \to B$ est un homomorphisme quasi-fini plat d'anneaux noethériens,
alors $\omega_{B/A}$ est un $B$-module inversible si et seulement si
$\omega_{B \otimes_A \kappa(\mathfrak p)/\kappa(\mathfrak p)}$
est un $B \otimes_A \kappa(\mathfrak p)$-module inversible
pour tout idéal premier $\mathfrak p \subset A$.
\item Si $Y \to X$ est un morphisme quasi-fini plat de
schémas noethériens, alors $\omega_{Y/X}$ est inversible
si et seulement si $\omega_{Y_x/x}$ est inversible pour tout $x \in X$.
\end{enumerate}
\end{lemma}

\begin{proof}
Preuve de (1). Comme $A \to B$ est plat, le module
$\omega_{B/A}$ est $A$-plat, voir le lemme \ref{lemma-dualizing-base-flat-flat}.
Ainsi $\omega_{B/A}$ est un $B$-module inversible si et seulement si
$\omega_{B/A} \otimes_A \kappa(\mathfrak p)$
est un $B \otimes_A \kappa(\mathfrak p)$-module inversible pour
tout idéal premier $\mathfrak p \subset A$, voir Compléments sur les morphismes, lemme
\ref{more-morphisms-lemma-flat-and-free-at-point-fibre}.
Toujours par platitude de $A \to B$, la
formation de $\omega_{B/A}$ commute au changement de base, voir le
lemme \ref{lemma-dualizing-base-change-of-flat}.
On voit donc que, sous l'hypothèse de platitude, l'inversibilité du module dualisant relatif
équivaut à l'inversibilité
du module dualisant relatif pour les homomorphismes
$\kappa(\mathfrak p) \to B \otimes_A \kappa(\mathfrak p)$.

\medskip\noindent
L'assertion (2) découle de (1) et de ce que, localement sur les ouverts affines,
les modules dualisants sont donnés par leurs analogues algébriques, voir la
remarque \ref{remark-relative-dualizing-for-quasi-finite}.
\end{proof}

\begin{lemma}
\label{lemma-dim-zero-global-complete-intersection-over-field}
Soit $k$ un corps. Soit $B = k[x_1, \ldots, x_n]/(f_1, \ldots, f_n)$
une intersection complète globale sur $k$, de dimension $0$.
Alors $\omega_{B/k}$ est inversible.
\end{lemma}

\begin{proof}
D'après le lemme de normalisation de Noether, voir
Algèbre, lemme \ref{algebra-lemma-Noether-normalization},
il existe une injection finie $k \to B$, c'est-à-dire
$\dim_k(B) < \infty$. Ainsi $\omega_{B/k} = \Hom_k(B, k)$
comme $B$-module.
D'après Complexes dualisants, lemme \ref{dualizing-lemma-dualizing-finite},
$R\Hom(B, k)$ est un complexe dualisant de $B$,
et Complexes dualisants, lemme \ref{dualizing-lemma-RHom-ext},
montre que $R\Hom(B, k)$ est égal à $\omega_{B/k}$
placé en degré $0$. Il suffit donc de montrer que
$B$ est de Gorenstein
(Complexes dualisants, lemme \ref{dualizing-lemma-gorenstein}).
C'est le cas d'après Complexes dualisants, lemme
\ref{dualizing-lemma-gorenstein-lci}.
\end{proof}

\begin{lemma}
\label{lemma-dualizing-syntomic-quasi-finite}
Soit $f : Y \to X$ un morphisme de schémas localement noethériens. Si $f$
vérifie les conditions équivalentes du lemme \ref{lemma-syntomic-quasi-finite},
alors $\omega_{Y/X}$ est un $\mathcal{O}_Y$-module inversible.
\end{lemma}

\begin{proof}
On peut supposer que $A \to B$ est une intersection complète globale
relative de la forme $B = A[x_1, \ldots, x_n]/(f_1, \ldots, f_n)$,
et il faut montrer que $\omega_{B/A}$ est inversible.
Cela résulte de la combinaison des lemmes \ref{lemma-characterize-invertible} et
\ref{lemma-dim-zero-global-complete-intersection-over-field}.
\end{proof}

\begin{example}
\label{example-universal-quasi-finite-syntomic}
Soient $n \geq 1$ et $d \geq 1$ des entiers. Soit $T$ l'ensemble des
multi-indices $E = (e_1, \ldots, e_n)$ tels que $e_i \geq 0$ et
$\sum e_i \leq d$. Considérons l'anneau
$$
A = \mathbf{Z}[a_{i, E} ; 1 \leq i \leq n, E \in T]
$$
Dans $A[x_1, \ldots, x_n]$, considérons les éléments
$f_i = \sum_{E \in T} a_{i, E} x^E$, où $x^E = x_1^{e_1} \ldots x_n^{e_n}$
suivant la notation usuelle. Considérons la $A$-algèbre
$$
B = A[x_1, \ldots, x_n]/(f_1, \ldots, f_n)
$$
Notons $X_{n, d} = \Spec(A)$ et soit $Y_{n, d} \subset \Spec(B)$
le plus grand sous-schéma ouvert sur lequel la restriction du
morphisme $\Spec(B) \to \Spec(A) = X_{n, d}$ est quasi-finie, voir
Algèbre, lemme \ref{algebra-lemma-quasi-finite-open}.
\end{example}

\begin{lemma}
\label{lemma-universal-quasi-finite-syntomic-etale}
Avec les notations de l'exemple \ref{example-universal-quasi-finite-syntomic},
les schémas $X_{n, d}$ et $Y_{n, d}$ sont réguliers et irréductibles,
le morphisme $Y_{n, d} \to X_{n, d}$ est localement quasi-fini et
syntomique, et il existe un sous-schéma ouvert dense $V \subset Y_{n, d}$
tel que $Y_{n, d} \to X_{n, d}$ se restreigne en un morphisme étale
$V \to X_{n, d}$.
\end{lemma}

\begin{proof}
Le schéma $X_{n, d}$ est le spectre de l'anneau de polynômes $A$.
Ainsi $X_{n, d}$ est régulier et irréductible. Puisque l'on peut écrire
$$
f_i = a_{i, (0, \ldots, 0)} +
\sum\nolimits_{E \in T, E \not = (0, \ldots, 0)} a_{i, E} x^E
$$
on voit que l'anneau $B$ est isomorphe à l'anneau de polynômes
en $x_1, \ldots, x_n$ et en les éléments $a_{i, E}$ tels que
$E \not = (0, \ldots, 0)$. Ainsi $\Spec(B)$ est un schéma irréductible et
régulier, de même que l'ouvert $Y_{n, d}$. Le morphisme
$Y_{n, d} \to X_{n, d}$ est localement quasi-fini et syntomique d'après le
lemme \ref{lemma-syntomic-quasi-finite}. Pour trouver $V$, il suffit
de trouver un point où $Y_{n, d} \to X_{n, d}$ est étale
(le lieu des points où un morphisme est étale est ouvert par
définition). Il suffit donc de trouver un point de $X_{n, d}$
où la fibre de $Y_{n, d} \to X_{n, d}$ est non vide et étale, voir
Morphismes, lemme \ref{morphisms-lemma-etale-at-point}. Choisissons
le point correspondant à l'homomorphisme d'anneaux $\chi : A \to \mathbf{Q}$
envoyant $f_i$ sur $x_i^d - 1$. Alors
$$
B \otimes_{A, \chi} \mathbf{Q} =
\mathbf{Q}[x_1, \ldots, x_n]/(x_1^d - 1, \ldots, x_n^d - 1)
$$
est une algèbre étale non nulle sur $\mathbf{Q}$.
\end{proof}

\begin{lemma}
\label{lemma-locally-comes-from-universal}
Soit $f : Y \to X$ un morphisme de schémas. Si $f$ vérifie les conditions équivalentes
du lemme \ref{lemma-syntomic-quasi-finite}, alors, pour tout
$y \in Y$, il existe $n, d$ et un diagramme commutatif
$$
\xymatrix{
Y \ar[d] &
V \ar[d] \ar[l] \ar[r] &
Y_{n, d} \ar[d] \\
X & U \ar[l] \ar[r] &
X_{n, d}
}
$$
où $U \subset X$ et $V \subset Y$ sont ouverts, avec $y \in V$,
où $Y_{n, d} \to X_{n, d}$
est comme dans l'exemple \ref{example-universal-quasi-finite-syntomic}, et
où le carré de droite est cartésien.
\end{lemma}

\begin{proof}
D'après le lemme \ref{lemma-syntomic-quasi-finite},
on peut choisir $U$ et $V$ affines, de sorte que
$U = \Spec(R)$ et $V = \Spec(S)$ avec
$S = R[y_1, \ldots, y_n]/(g_1, \ldots, g_n)$.
Avec les notations de l'exemple \ref{example-universal-quasi-finite-syntomic},
si l'on choisit $d$ assez grand, on peut écrire chaque $g_i$ sous la forme
$g_i = \sum_{E \in T} g_{i, E}y^E$ avec $g_{i, E} \in R$.
L'application $A \to R$ envoyant $a_{i, E}$ sur $g_{i, E}$
et l'application $B \to S$ envoyant $x_i \to y_i$ donnent alors un diagramme cocartésien
d'anneaux
$$
\xymatrix{
S & B \ar[l] \\
R \ar[u] & A \ar[l] \ar[u]
}
$$
ce qui démontre le lemme.
\end{proof}











\section{Morphismes finis syntomiques}
\label{section-finite-syntomic}

\noindent
Cette section est l'analogue de la section \ref{section-quasi-finite-syntomic}
pour les morphismes finis syntomiques.

\begin{lemma}
\label{lemma-syntomic-finite}
Soit $f : Y \to X$ un morphisme de schémas. Les conditions suivantes sont équivalentes :
\begin{enumerate}
\item $f$ est fini et syntomique ;
\item $f$ est fini, plat et localement d'intersection complète ;
\item $f$ est fini, plat, localement de présentation finie,
et les fibres de $f$ sont localement des intersections complètes ;
\item $f$ est fini et, pour tout $x \in X$, il existe un
ouvert affine $x \in U = \Spec(A) \subset X$, un entier $n$
et $f_1, \ldots, f_n \in A[x_1, \ldots, x_n]$ tels que
$f^{-1}(U)$ soit isomorphe au spectre de
$A[x_1, \ldots, x_n]/(f_1, \ldots, f_n)$ ;
\item {\emergencystretch=1em $f$ est fini, plat, localement de présentation finie,
et $\NL_{X/Y}$ est de Tor-amplitude contenue dans $[-1, 0]$ ;\par}
\item $f$ est fini, plat, localement de présentation finie, et
$\NL_{X/Y}$ est parfait de rang $0$ et de Tor-amplitude contenue dans $[-1, 0]$.
\end{enumerate}
\end{lemma}

\begin{proof}
L'équivalence de (1), (2), (3), (5) et (6)
ainsi que l'implication (4) $\Rightarrow$ (1) résultent immédiatement
du lemme \ref{lemma-syntomic-quasi-finite}. Supposons les conditions équivalentes
(1), (2), (3), (5), (6) satisfaites.
Choisissons un point $x \in X$ et un voisinage ouvert affine $U = \Spec(A)$
de $x$ dans $X$, et supposons que $x$ corresponde à l'idéal premier
$\mathfrak p \subset A$. Écrivons $f^{-1}(U) = \Spec(B)$.
Écrivons $B = A[x_1, \ldots, x_n]/I$. Puisque $\NL_{B/A}$
est parfait de Tor-amplitude contenue dans $[-1, 0]$ d'après (6),
on voit que $I/I^2$ est un $B$-module fini localement libre
de rang $n$. Puisque $B_\mathfrak p$ est semi-local,
$(I/I^2)_\mathfrak p$ est libre de rang $n$, voir
Algèbre, lemme \ref{algebra-lemma-locally-free-semi-local-free}.
Ainsi, en remplaçant $A$ par sa localisation principale en
un élément n'appartenant pas à $\mathfrak p$, on peut supposer que $I/I^2$
est un $B$-module libre de rang $n$.
D'après Algèbre, lemme \ref{algebra-lemma-huber},
on peut donc trouver une présentation de $B$ sur $A$
ayant autant de variables que d'équations. Autrement dit,
on peut supposer $B = A[x_1, \ldots, x_n]/(f_1, \ldots, f_n)$.
Cela démontre (4).
\end{proof}

\begin{example}
\label{example-universal-finite-syntomic}
Soit $d \geq 1$ un entier. Considérons des variables
$a_{ij}^l$ pour $1 \leq i, j, l \leq d$, et notons
$$
A_d = \mathbf{Z}[a_{ij}^k]/J
$$
où $J$ est l'idéal engendré par les éléments
$$
\left\{
\begin{matrix}
\sum_l a_{ij}^la_{lk}^m - \sum_l a_{il}^ma_{jk}^l & \forall i, j, k, m \\
a_{ij}^k - a_{ji}^k & \forall i, j, k \\
a_{i1}^j - \delta_{ij} & \forall i, j
\end{matrix}
\right.
$$
où $\delta_{ij}$ désigne le symbole de Kronecker.
On définit une $A_d$-algèbre $B_d$ comme suit : comme $A_d$-module, on pose
$$
B_d = A_d e_1 \oplus \ldots \oplus A_d e_d
$$
La structure d'algèbre est donnée par $A_d \to B_d$ envoyant $1$ sur $e_1$.
La multiplication sur $B_d$ est l'application $A_d$-bilinéaire
$$
m : B_d \times B_d \longrightarrow B_d, \quad
m(e_i, e_j) = \sum a_{ij}^k e_k
$$
On vérifie sans difficulté que les relations ci-dessus
imposent exactement une structure de $A_d$-algèbre.
Le morphisme
$$
\pi_d : Y_d = \Spec(B_d) \longrightarrow \Spec(A_d) = X_d
$$
est le morphisme fini libre « universel » de rang $d$.
\end{example}

\begin{lemma}
\label{lemma-universal-finite-syntomic}
Avec les notations de l'exemple \ref{example-universal-finite-syntomic},
il existe un sous-schéma ouvert $U_d \subset X_d$ possédant la propriété suivante :
un morphisme de schémas $X \to X_d$ se factorise par $U_d$ si et seulement
si $Y_d \times_{X_d} X \to X$ est syntomique.
\end{lemma}

\begin{proof}
Rappelons qu'être syntomique équivaut à être plat et
localement d'intersection complète, voir
Compléments sur les morphismes, lemme \ref{more-morphisms-lemma-flat-lci}.
L'ensemble $W_d \subset Y_d$ des points où $\pi_d$ est de Koszul
est ouvert dans $Y_d$, et sa formation commute à tout changement de base, voir
Compléments sur les morphismes, lemme \ref{more-morphisms-lemma-base-change-lci-fibres}.
Puisque $\pi_d$ est fini, donc fermé,
$Z = \pi_d(Y_d \setminus W_d)$ est fermé. Comme on a manifestement $U_d = X_d \setminus Z$
et que sa formation commute au changement de base, le lemme
en résulte.
\end{proof}

\begin{lemma}
\label{lemma-universal-finite-syntomic-smooth}
Avec les notations de l'exemple \ref{example-universal-finite-syntomic}
et $U_d$ comme dans le lemme \ref{lemma-universal-finite-syntomic},
le schéma $U_d$ est lisse sur $\Spec(\mathbf{Z})$.
\end{lemma}

\begin{proof}
Utilisons Compléments sur les morphismes, lemme
\ref{more-morphisms-lemma-lifting-along-artinian-at-point},
pour montrer que $U_d \to \Spec(\mathbf{Z})$ est lisse.
Supposons donc que $\Spec(A) \to U_d$ soit un morphisme
et que $A' \to A$ soit une petite extension. Alors $B = A \otimes_{A_d} B_d$
est une $A$-algèbre finie libre, syntomique sur $A$
(par construction de $U_d$). D'après
Lissage des homomorphismes d'anneaux, proposition \ref{smoothing-proposition-lift-smooth},
il existe un homomorphisme syntomique d'anneaux $A' \to B'$ tel que
$B \cong B' \otimes_{A'} A$. Posons $e'_1 = 1 \in B'$. Pour $1 < i \leq d$,
choisissons des relèvements $e'_i \in B'$ des éléments
$1 \otimes e_i \in A \otimes_{A_d} B_d = B$. Alors $e'_1, \ldots, e'_d$
est une base de $B'$ sur $A'$ (voir par exemple Algèbre, lemme
\ref{algebra-lemma-local-artinian-basis-when-flat}).
On peut donc écrire $e'_i e'_j = \sum \alpha_{ij}^l e'_l$ pour d'uniques
éléments $\alpha_{ij}^l \in A'$ vérifiant les relations
$\sum_l \alpha_{ij}^l \alpha_{lk}^m = \sum_l \alpha_{il}^m \alpha _{jk}^l$
et $\alpha_{ij}^k = \alpha_{ji}^k$ et $\alpha_{i1}^j = \delta_{ij}$
dans $A'$. Cela détermine un morphisme $\Spec(A') \to X_d$ en
envoyant $a_{ij}^l \in A_d$ sur $\alpha_{ij}^l \in A'$. Ce morphisme
coïncide avec le morphisme donné $\Spec(A) \to U_d$. Puisque $\Spec(A')$
et $\Spec(A)$ ont le même espace topologique sous-jacent, on obtient
le relèvement voulu $\Spec(A') \to U_d$, et l'on
conclut que $U_d$ est lisse sur $\mathbf{Z}$.
\end{proof}

\begin{lemma}
\label{lemma-universal-finite-syntomic-etale}
Avec les notations de l'exemple \ref{example-universal-finite-syntomic},
considérons le sous-schéma ouvert $U'_d \subset X_d$ au-dessus duquel
$\pi_d$ est étale. Alors $U'_d$ est une partie dense de
l'ouvert $U_d$ du lemme \ref{lemma-universal-finite-syntomic}.
\end{lemma}

\begin{proof}
En raisonnant exactement comme dans la démonstration du
lemme \ref{lemma-universal-finite-syntomic}, à l'aide de
Morphismes, lemme \ref{morphisms-lemma-set-points-where-fibres-etale},
on obtient un plus grand ouvert $U'_d \subset X_d$ au-dessus duquel $\pi_d$ est
étale. De plus, un morphisme étale étant syntomique, on voit
que $U'_d \subset U_d$. Pour achever la démonstration, il faut montrer
que $U'_d \subset U_d$ est dense. Soit $u : \Spec(k) \to U_d$ un morphisme,
où $k$ est un corps. Soit $B = k \otimes_{A_d} B_d$ comme dans la
démonstration du lemme \ref{lemma-universal-finite-syntomic-smooth}.
Nous allons montrer qu'il existe un anneau local intègre $A'$ de corps résiduel $k$
et une $A'$-algèbre finie syntomique $B'$ avec $B = k \otimes_{A'} B'$
dont la fibre générique est étale. Exactement comme au paragraphe précédent,
cela déterminera un morphisme $\Spec(A') \to U_d$ envoyant le
point générique dans $U'_d$ et le point fermé sur $u$, ce qui
achèvera la démonstration.

\medskip\noindent
D'après le lemme \ref{lemma-syntomic-finite}, (4), on peut choisir une présentation\newline
$B = k[x_1, \ldots, x_n]/(f_1, \ldots, f_n)$.\newline
Soit $d'$ le maximum des degrés totaux des polynômes $f_1, \ldots, f_n$.
Soit $Y_{n, d'} \to X_{n, d'}$ comme dans
l'exemple \ref{example-universal-quasi-finite-syntomic}.
Par construction, il existe un morphisme $u' : \Spec(k) \to X_{n, d'}$
tel que
$$
\Spec(B) \cong Y_{n, d'} \times_{X_{n, d'}, u'} \Spec(k)
$$
Notons $A = \mathcal{O}_{X_{n, d'}, u'}^h$ le hensélisé de
l'anneau local de $X_{n, d'}$ à l'image de $u'$. On peut alors écrire
$$
Y_{n, d'} \times_{X_{n, d'}} \Spec(A) = Z \amalg W
$$
avec $Z \to \Spec(A)$ fini et $W \to \Spec(A)$ de fibre
fermée vide, voir
Algèbre, lemme \ref{algebra-lemma-characterize-henselian}, (13),
ou la discussion dans Compléments sur les morphismes, section
\ref{more-morphisms-section-etale-localization}.
D'après le lemme \ref{lemma-universal-quasi-finite-syntomic-etale},
l'anneau local $A$ est régulier (on utilise également ici
Compléments d'algèbre, lemme \ref{more-algebra-lemma-henselization-regular}),
et le morphisme $Z \to \Spec(A)$ est étale au-dessus du point générique de
$\Spec(A)$ (car celui-ci s'envoie sur le point générique de $X_{d, n'}$).
Par construction, $Z \times_{\Spec(A)} \Spec(k) \cong \Spec(B)$.
Cela donne le résultat voulu, à ceci près que l'application du
corps résiduel de $A$ vers $k$ n'est pas nécessairement un isomorphisme.
D'après Algèbre, lemme \ref{algebra-lemma-flat-local-given-residue-field},
il existe un homomorphisme local plat d'anneaux $A \to A'$ tel que le corps résiduel
de $A'$ soit $k$. Si $A'$ n'est pas intègre, on choisit un
idéal premier minimal $\mathfrak p \subset A'$ (situé au-dessus de
l'unique idéal premier minimal de $A$ par platitude) et l'on remplace
$A'$ par $A'/\mathfrak p$. On prend pour $B'$ l'unique $A'$-algèbre
telle que $Z \times_{\Spec(A)} \Spec(A') = \Spec(B')$.
La démonstration est achevée.
\end{proof}

\begin{remark}
\label{remark-universal-finite-syntomic-smooth-top}
Soit $\pi_d : Y_d \to X_d$ comme dans
l'exemple \ref{example-universal-finite-syntomic}.
Soit $U_d \subset X_d$ le plus grand ouvert au-dessus duquel
$V_d = \pi_d^{-1}(U_d)$ est fini syntomique, comme dans le
lemme \ref{lemma-universal-finite-syntomic}.
Alors $V_d$ est également lisse sur $\mathbf{Z}$.
(Bien entendu, le morphisme $V_d \to U_d$ n'est pas lisse lorsque $d \geq 2$.)
En raisonnant comme dans la démonstration du lemme \ref{lemma-universal-finite-syntomic-smooth},
on voit que cela correspond au problème de déformation
suivant : étant donné une petite extension $C' \to C$ et
une $C$-algèbre finie syntomique $B$ munie d'une section $B \to C$,
trouver une $C'$-algèbre finie syntomique $B'$ et une section $B' \to C'$
dont le produit tensoriel avec $C$ redonne $B \to C$.
D'après le lemme \ref{lemma-syntomic-finite}, on peut écrire
$B = C[x_1, \ldots, x_n]/(f_1, \ldots, f_n)$ comme
une intersection complète globale relative.
Après un changement de coordonnées, on peut supposer
$x_1, \ldots, x_n$ dans le noyau de $B \to C$.
Les polynômes $f_i$ ont alors un terme constant nul.
Choisissons des relèvements quelconques $f'_i \in C'[x_1, \ldots, x_n]$ de $f_i$
ayant un terme constant nul. Alors
$B' = C'[x_1, \ldots, x_n]/(f'_1, \ldots, f'_n)$,
avec la section $B' \to C'$ envoyant $x_i$ sur zéro, convient.
\end{remark}

\begin{lemma}
\label{lemma-locally-comes-from-universal-finite}
Soit $f : Y \to X$ un morphisme de schémas. Si $f$ vérifie les conditions équivalentes
du lemme \ref{lemma-syntomic-finite}, alors, pour tout
$x \in X$, il existe un $d$ et un diagramme commutatif
$$
\xymatrix{
Y \ar[d] &
V \ar[d] \ar[l] \ar[r] &
V_d \ar[d] \ar[r] &
Y_d \ar[d]^{\pi_d}\\
X &
U \ar[l] \ar[r] &
U_d \ar[r] &
X_d
}
$$
possédant les propriétés suivantes :
\begin{enumerate}
\item $U \subset X$ est ouvert, $x \in U$ et $V = f^{-1}(U)$ ;
\item $\pi_d : Y_d \to X_d$ est comme dans
l'exemple \ref{example-universal-finite-syntomic} ;
\item $U_d \subset X_d$ est comme dans le lemme \ref{lemma-universal-finite-syntomic}
et $V_d = \pi_d^{-1}(U_d) \subset Y_d$ ;
\item le carré du milieu est cartésien.
\end{enumerate}
\end{lemma}

\begin{proof}
Choisissons un voisinage ouvert affine $U = \Spec(A) \subset X$ de $x$.
Écrivons $V = f^{-1}(U) = \Spec(B)$. Alors $B$ est un
$A$-module fini localement libre, et l'inclusion $A \subset B$ est localement un facteur direct.
En rétrécissant $U$, on peut donc choisir une base $1 = e_1, e_2, \ldots, e_d$
de $B$ comme $A$-module. Écrivons
$e_i e_j = \sum \alpha_{ij}^l e_l$ pour d'uniques
éléments $\alpha_{ij}^l \in A$ vérifiant les relations
$\sum_l \alpha_{ij}^l \alpha_{lk}^m = \sum_l \alpha_{il}^m \alpha _{jk}^l$
et $\alpha_{ij}^k = \alpha_{ji}^k$ et $\alpha_{i1}^j = \delta_{ij}$
dans $A$. Cela détermine un morphisme $\Spec(A) \to X_d$ en envoyant
$a_{ij}^l \in A_d$ sur $\alpha_{ij}^l \in A$. Par construction,
$V \cong \Spec(A) \times_{X_d} Y_d$. La définition de $U_d$
montre que $\Spec(A) \to X_d$ se factorise par $U_d$.
La démonstration est achevée.
\end{proof}










\section{Une formule pour la différente}
\label{section-formula-different}

\noindent
Dans cette section, nous présentons les résultats dus à Tate dans \cite[Appendix A]{Mazur-Roberts}.
Dans notre terminologie, ils montreront que la différente est
égale à la différente de Kähler pour un morphisme plat, quasi-fini,
localement d'intersection complète.
Calculons d'abord la différente de Noether dans un cas particulier.

\begin{lemma}
\label{lemma-tate}
\begin{reference}
\cite[Appendix]{Mazur-Roberts}
\end{reference}
Soit $A \to P$ un homomorphisme d'anneaux. Soit $f_1, \ldots, f_n \in P$ une
suite régulière au sens de Koszul. Supposons $B = P/(f_1, \ldots, f_n)$
plat sur $A$. Soit $g_1, \ldots, g_n \in P \otimes_A B$
une suite régulière au sens de Koszul engendrant le noyau de l'application de multiplication
$P \otimes_A B \to B$. Écrivons $f_i \otimes 1 = \sum g_{ij} g_j$.
Alors l'annulateur de $\Ker(B \otimes_A B \to B)$ est un idéal principal
engendré par l'image de $\det(g_{ij})$.
\end{lemma}

\begin{proof}
Le complexe de Koszul $K_\bullet = K(P, f_1, \ldots, f_n)$ est une résolution
de $B$ par des $P$-modules libres de type fini. Le complexe de Koszul
$M_\bullet = K(P \otimes_A B, g_1, \ldots, g_n)$ est une résolution
de $B$ par des $P \otimes_A B$-modules libres de type fini. Il existe un morphisme de
complexes
$$
K_\bullet \longrightarrow M_\bullet
$$
donné en degré $1$ par la matrice $(g_{ij})$ et
en degré $n$ par $\det(g_{ij})$. Voir
Compléments d'algèbre, lemme \ref{more-algebra-lemma-functorial}.
Comme $B$ est un $A$-module plat, on peut considérer $M_\bullet$ comme un complexe
de $P$-modules plats (via $P \to P \otimes_A B$, $p \mapsto p \otimes 1$).
On peut donc utiliser les deux complexes pour calculer $\text{Tor}_*^P(B, B)$,
et l'application affichée devient un quasi-isomorphisme après produit tensoriel
avec $B$. Il est clair que $H_n(K_\bullet \otimes_P B) = B$.
D'autre part, $H_n(M_\bullet \otimes_P B)$ est le noyau de
$$
B \otimes_A B \xrightarrow{g_1, \ldots, g_n} (B \otimes_A B)^{\oplus n}
$$
Puisque $g_1, \ldots, g_n$ engendrent le noyau de $B \otimes_A B \to B$,
le lemme est démontré.
\end{proof}

\begin{lemma}
\label{lemma-quasi-finite-complete-intersection}
Soient $A$ un anneau, $n \geq 1$ et
$h, f_1, \ldots, f_n \in A[x_1, \ldots, x_n]$.
Posons $B = A[x_1, \ldots, x_n, 1/h]/(f_1, \ldots, f_n)$.
Supposons que $B$ soit quasi-fini sur $A$.
Alors
\begin{enumerate}
\item $B$ est plat sur $A$ et $A \to B$ est une intersection complète locale
relative ;
\item l'annulateur $J$ de $I = \Ker(B \otimes_A B \to B)$
est libre de rang $1$ sur $B$ ;
\item la différente de Noether de $B$ sur $A$ est engendrée
par $\det(\partial f_i/\partial x_j)$ dans $B$.
\end{enumerate}
\end{lemma}

\begin{proof}
Remarquons que
$B = A[x, x_1, \ldots, x_n]/(xh - 1, f_1, \ldots, f_n)$
est une intersection complète globale relative sur $A$, voir
Algèbre, définition
\ref{algebra-definition-relative-global-complete-intersection}.
D'après Algèbre, lemme \ref{algebra-lemma-relative-global-complete-intersection},
on voit que $B$ est plat sur $A$.

\medskip\noindent
Posons $P' = A[x, x_1, \ldots, x_n]$ et
$P = P'/(xh - 1) = A[x_1, \ldots, x_n, 1/h]$.
On a alors $P' \to P \to B$.
D'après Compléments d'algèbre, lemme
\ref{more-algebra-lemma-relative-global-complete-intersection-koszul},
la suite $xh - 1, f_1, \ldots, f_n$ est régulière au sens de Koszul
dans $P'$. Puisque $xh - 1$ est une suite régulière au sens de Koszul de longueur
un dans $P'$ (par exemple d'après le même lemme), on en déduit que
$f_1, \ldots, f_n$ est une suite régulière au sens de Koszul dans $P$ d'après
Compléments d'algèbre, lemme \ref{more-algebra-lemma-truncate-koszul-regular}.

\medskip\noindent
Soit $g_i \in P \otimes_A B$ l'image de $x_i \otimes 1 - 1 \otimes x_i$.
Utilisons les notations abrégées $y_i = x_i \otimes 1$ et $z_i = 1 \otimes x_i$
dans $A[x_1, \ldots, x_n] \otimes_A A[x_1, \ldots, x_n]$,
de sorte que $g_i$ soit l'image de $y_i - z_i$. Pour un polynôme
$f \in A[x_1, \ldots, x_n]$, on écrit $f(y) = f \otimes 1$
et $f(z) = 1 \otimes f$ dans le produit tensoriel ci-dessus.
On a alors
$$
P \otimes_A B/(g_1, \ldots, g_n) =
\frac{A[y_1, \ldots, y_n, z_1, \ldots, z_n, \frac{1}{h(y)h(z)}]}
{(f_1(z), \ldots, f_n(z), y_1 - z_1, \ldots, y_n - z_n)}
$$
qui est manifestement isomorphe à $B$. Par les mêmes arguments
que ci-dessus, on trouve donc que $f_1(z), \ldots, f_n(z), y_1 - z_1, \ldots, y_n - z_n$
est une suite régulière au sens de Koszul dans
$A[y_1, \ldots, y_n, z_1, \ldots, z_n, \frac{1}{h(y)h(z)}]$.
La suite $f_1(z), \ldots, f_n(z)$ est régulière au sens de Koszul dans
$A[y_1, \ldots, y_n, z_1, \ldots, z_n, \frac{1}{h(y)h(z)}]$
par platitude de l'application
$$
P \longrightarrow A[y_1, \ldots, y_n, z_1, \ldots, z_n,
\textstyle{\frac{1}{h(y)h(z)}}],\quad x_i \longmapsto z_i
$$
et d'après Compléments d'algèbre, lemme
\ref{more-algebra-lemma-koszul-regular-flat-base-change}.
D'après Compléments d'algèbre, lemme \ref{more-algebra-lemma-truncate-koszul-regular},
on conclut que $g_1, \ldots, g_n$ est une suite régulière
dans $P \otimes_A B$.

\medskip\noindent
Nous avons maintenant vérifié toutes les hypothèses du lemme \ref{lemma-tate}
ci-dessus, avec $P$, $f_1, \ldots, f_n$ et $g_i \in P \otimes_A B$ comme ci-dessus.
En particulier, l'annulateur $J$ de $I$ est libre, engendré par un
élément $\delta$ sur $B$.
Posons $f_{ij} = \partial f_i/\partial x_j \in A[x_1, \ldots, x_n]$.
Un calcul élémentaire montre que l'on peut écrire
$$
f_i(y) =
f_i(z_1 + g_1, \ldots, z_n + g_n) =
f_i(z) + \sum\nolimits_j f_{ij}(z) g_j +
\sum\nolimits_{j, j'} F_{ijj'}g_jg_{j'}
$$
avec $F_{ijj'} \in A[y_1, \ldots, y_n, z_1, \ldots, z_n]$.
En passant à l'image dans $P \otimes_A B$, les termes $f_i(z)$ deviennent
nuls et l'on obtient
$$
f_i \otimes 1 = \sum\nolimits_j
\left(1 \otimes f_{ij} + \sum\nolimits_{j'} F_{ijj'}g_{j'}\right)g_j
$$
On conclut donc du lemme \ref{lemma-tate}
que $\delta = \det(g_{ij})$, avec
$g_{ij} = 1 \otimes f_{ij} + \sum_{j'} F_{ijj'}g_{j'}$.
Puisque $g_{j'}$ a une image nulle dans $B$, on en déduit
que l'image de $\det(\partial f_i/\partial x_j)$ dans $B$
engendre la différente de Noether de $B$ sur $A$.
\end{proof}

\begin{lemma}
\label{lemma-different-syntomic-quasi-finite}
Soit $f : Y \to X$ un morphisme de schémas noethériens. Si $f$
vérifie les conditions équivalentes du lemme \ref{lemma-syntomic-quasi-finite},
alors la différente $\mathfrak{D}_f$ de $f$ est la différente de Kähler
de $f$.
\end{lemma}

\begin{proof}
D'après les lemmes \ref{lemma-flat-gorenstein-agree-noether} et
\ref{lemma-dualizing-syntomic-quasi-finite},
la différente de $f$ coïncide, localement sur les ouverts affines, avec la
différente de Noether. Le lemme résulte alors du
calcul de la différente de Noether et de la différente de Kähler
sur les morceaux affines standards, effectué dans les
lemmes \ref{lemma-kahler-different-complete-intersection} et
\ref{lemma-quasi-finite-complete-intersection}.
\end{proof}

\begin{lemma}
\label{lemma-different-quasi-finite-complete-intersection}
Soient $A$ un anneau, $n \geq 1$ et
$h, f_1, \ldots, f_n \in A[x_1, \ldots, x_n]$.
Posons $B = A[x_1, \ldots, x_n, 1/h]/(f_1, \ldots, f_n)$.
Supposons que $B$ soit quasi-fini sur $A$.
Il existe alors un isomorphisme $B \to \omega_{B/A}$
envoyant $\det(\partial f_i/\partial x_j)$ sur $\tau_{B/A}$.
\end{lemma}

\begin{proof}
Soit $J$ l'annulateur de $\Ker(B \otimes_A B \to B)$.
D'après le lemme \ref{lemma-quasi-finite-complete-intersection},
l'homomorphisme $A \to B$ est plat et
$J$ est un $B$-module libre ayant un générateur $\xi$ dont l'image est
$\det(\partial f_i/\partial x_j)$ dans $B$.
Le lemme découle donc du
lemme \ref{lemma-noether-different-flat-quasi-finite}
et du fait (lemme \ref{lemma-dualizing-syntomic-quasi-finite})
que $\omega_{B/A}$ est un $B$-module inversible.
(Attention : il est nécessaire de montrer que $\omega_{B/A}$
est inversible, car un $B$-module de type fini $M$ tel
que $\Hom_B(M, B) \cong B$ n'est pas nécessairement libre.)
\end{proof}

\begin{example}
\label{example-different-for-monogenic}
Soit $A$ un anneau noethérien. Soient $f, h \in A[x]$ tels que
$$
B = (A[x]/(f))_h = A[x, 1/h]/(f)
$$
soit quasi-fini sur $A$. Soit $f' \in A[x]$ la dérivée
de $f$ par rapport à $x$. L'idéal $\mathfrak{D} = (f') \subset B$
est la différente de Noether de $B$ sur $A$,
est la différente de Kähler de $B$ sur $A$, et
est l'idéal dont le faisceau quasi-cohérent d'idéaux associé est la
différente de $\Spec(B)$ sur $\Spec(A)$.
\end{example}

\begin{lemma}
\label{lemma-discriminant-quasi-finite-morphism-smooth}
Soit $S$ un schéma noethérien. Soient $X$, $Y$ des schémas lisses
de dimension relative $n$ sur $S$. Soit $f : Y \to X$ un
morphisme localement quasi-fini sur $S$.
Alors $f$ est plat et le sous-schéma fermé $R \subset Y$
défini par la différente de $f$ est le sous-schéma fermé localement principal
défini par
$$
\wedge^n(\text{d}f) \in
\Gamma(Y,
(f^*\Omega^n_{X/S})^{\otimes -1} \otimes_{\mathcal{O}_Y} \Omega^n_{Y/S})
$$
Si $f$ est étale aux points associés de $Y$, alors $R$ est un
diviseur de Cartier effectif, et
$$
f^*\Omega^n_{X/S} \otimes_{\mathcal{O}_Y} \mathcal{O}(R) =
\Omega^n_{Y/S}
$$
comme faisceaux inversibles sur $Y$.
\end{lemma}

\begin{proof}
Pour montrer que $f$ est plat, il suffit de montrer que $Y_s \to X_s$
est plat pour tout $s \in S$ (Compléments sur les morphismes, lemme
\ref{more-morphisms-lemma-morphism-between-flat-Noetherian}).
La platitude de $Y_s \to X_s$ découle de
Algèbre, lemme \ref{algebra-lemma-CM-over-regular-flat}.
D'après Compléments sur les morphismes, lemme
\ref{more-morphisms-lemma-lci-permanence},
le morphisme $f$ est localement d'intersection complète.
L'assertion sur la différente découle donc de
l'assertion correspondante sur la différente de Kähler d'après le
lemme \ref{lemma-different-syntomic-quasi-finite}.
Enfin, puisque l'on a la suite exacte
$$
f^*\Omega_{X/S} \xrightarrow{\text{d}f} \Omega_{Y/S} \to \Omega_{Y/X} \to 0
$$
d'après Morphismes, lemme \ref{morphisms-lemma-triangle-differentials},
et que $\Omega_{X/S}$ et $\Omega_{Y/S}$ sont finis localement libres
de rang $n$ (Morphismes, lemme
\ref{morphisms-lemma-smooth-omega-finite-locally-free}),
l'assertion concernant la différente de Kähler résulte immédiatement de la définition
de l'idéal de Fitting d'ordre zéro. Si $f$ est étale aux points associés
de $Y$, alors $\wedge^n\text{d}f$ ne s'annule pas aux
points associés de $Y$, ce qui entraîne que l'équation locale
de $R$ est non diviseur de zéro. Ainsi $R$ est un diviseur de Cartier effectif.
L'isomorphisme canonique envoie $1$ sur $\wedge^n\text{d}f$, voir
Diviseurs, lemme \ref{divisors-lemma-characterize-OD}.
\end{proof}






\section{L'application de Tate}
\label{section-tate-map}

\noindent
Dans cette section, nous construisons un isomorphisme entre
le déterminant du complexe cotangent relatif et
le module dualisant relatif pour un morphisme localement quasi-fini
syntomique de schémas localement noethériens. Suivant
\cite[1.4.4]{Garel}, nous l'appelons l'application de Tate.
Notre méthode consiste à éviter autant que possible
les calculs locaux.

\medskip\noindent
Soit $Y \to X$ un morphisme localement quasi-fini syntomique de schémas.
Nous utiliserons dans cette section, sans le mentionner à nouveau, toutes les conditions
équivalentes données dans le lemme \ref{lemma-syntomic-quasi-finite}
pour cette notion. On voit en particulier que $\NL_{Y/X}$ est un objet parfait
de $D(\mathcal{O}_Y)$ de Tor-amplitude contenue dans $[-1, 0]$. On dispose donc
d'un module inversible canonique
$\det(\NL_{Y/X})$ sur $Y$ et d'une section globale
$$
\delta(\NL_{Y/X}) \in \Gamma(Y, \det(\NL_{Y/X}))
$$
Voir Catégories dérivées de schémas, lemme
\ref{perfect-lemma-determinant-two-term-complexes}.
Supposons donné un diagramme commutatif de schémas
$$
\xymatrix{
Y' \ar[r]_b \ar[d] & Y \ar[d] \\
X' \ar[r] & X
}
$$
dont les flèches verticales sont localement quasi-finies syntomiques et qui
induit un isomorphisme de $Y'$ avec un ouvert de $X' \times_X Y$.
L'application canonique
$$
Lb^*\NL_{Y/X} \longrightarrow \NL_{Y'/X'}
$$
est alors un quasi-isomorphisme d'après
Compléments sur les morphismes, lemme \ref{more-morphisms-lemma-base-change-NL-flat}.
On obtient donc un isomorphisme canonique
$b^*\det(\NL_{Y/X}) \to \det(\NL_{Y'/X'})$ qui envoie la
section canonique $\delta(\NL_{Y/X})$ sur $\delta(\NL_{Y'/ X'})$, voir
Catégories dérivées de schémas, remarque \ref{perfect-remark-functorial-det}.

\begin{remark}
\label{remark-local-description-delta}
Soit $Y \to X$ un morphisme localement quasi-fini syntomique de schémas.
Quelle est la forme locale du couple $(\det(\NL_{Y/X}), \delta(\NL_{Y/X}))$ ?
Choisissons des ouverts affines $V = \Spec(B) \subset Y$,
$U = \Spec(A) \subset X$ avec $f(V) \subset U$, un entier $n$ et
$f_1, \ldots, f_n \in A[x_1, \ldots, x_n]$ tels que
$B = A[x_1, \ldots, x_n]/(f_1, \ldots, f_n)$. Alors
$$
\NL_{B/A} = \left(
(f_1, \ldots, f_n)/(f_1, \ldots, f_n)^2
\longrightarrow
\bigoplus\nolimits_{i = 1, \ldots, n} B \text{d} x_i\right)
$$
et $(f_1, \ldots, f_n)/(f_1, \ldots, f_n)^2$ est libre de base
formée des classes $\overline{f}_i$. Voir la démonstration du
lemme \ref{lemma-syntomic-quasi-finite}.
Ainsi $\det(L_{B/A})$ est libre sur le générateur
$$
\text{d}x_1 \wedge \ldots \wedge \text{d}x_n
\otimes
(\overline{f}_1 \wedge \ldots \wedge \overline{f}_n)^{\otimes -1}
$$
et la section $\delta(\NL_{B/A})$ est l'élément
$$
\delta(\NL_{B/A}) =
\det(\partial f_j/ \partial x_i) \cdot
\text{d}x_1 \wedge \ldots \wedge \text{d}x_n
\otimes
(\overline{f}_1 \wedge \ldots \wedge \overline{f}_n)^{\otimes -1}
$$
par définition.
\end{remark}

\noindent
Soit $Y \to X$ un morphisme localement quasi-fini syntomique de
schémas localement noethériens. D'après les
remarques \ref{remark-relative-dualizing-for-quasi-finite} et
\ref{remark-relative-dualizing-for-flat-quasi-finite}, on dispose
d'un $\mathcal{O}_Y$-module cohérent $\omega_{Y/X}$ et d'une section globale
canonique
$$
\tau_{Y/X} \in \Gamma(Y, \omega_{Y/X})
$$
qui redonne, localement sur les ouverts affines, le couple $\omega_{B/A}, \tau_{B/A}$.
D'après le lemme \ref{lemma-dualizing-syntomic-quasi-finite}, le module
$\omega_{Y/X}$ est inversible. Supposons donné un diagramme commutatif de
schémas localement noethériens
$$
\xymatrix{
Y' \ar[r]_b \ar[d] & Y \ar[d] \\
X' \ar[r] & X
}
$$
dont les flèches verticales sont localement quasi-finies syntomiques et qui
induit un isomorphisme de $Y'$ avec un ouvert de $X' \times_X Y$.
Il existe alors une application canonique de changement de base
$$
b^*\omega_{Y/X} \longrightarrow \omega_{Y'/X'}
$$
qui est un isomorphisme
envoyant $\tau_{Y/X}$ sur $\tau_{Y'/X'}$. En effet, l'application de changement de base
dans le cadre affine est (\ref{equation-bc-dualizing}) ; c'est un
isomorphisme d'après le lemme \ref{lemma-dualizing-base-change-of-flat}, et elle
envoie $\tau_{Y/X}$ sur $\tau_{Y'/X'}$ d'après le
lemme \ref{lemma-trace-base-change}, (1).

\begin{proposition}
\label{proposition-tate-map}
Il existe une unique règle qui associe à tout morphisme localement quasi-fini syntomique
de schémas localement noethériens $Y \to X$ un isomorphisme
$$
c_{Y/X} : \det(\NL_{Y/X}) \longrightarrow \omega_{Y/X}
$$
vérifiant les deux propriétés suivantes :
\begin{enumerate}
\item la section $\delta(\NL_{Y/X})$ s'envoie sur $\tau_{Y/X}$ ;
\item la règle est compatible à la restriction aux ouverts et au
changement de base.
\end{enumerate}
\end{proposition}

\begin{proof}
Reformulons l'énoncé de la proposition. Considérons la catégorie
$\mathcal{C}$ dont les objets, notés $Y/X$, sont les morphismes localement quasi-finis syntomiques
$Y \to X$ de schémas localement noethériens, et dont les morphismes
$b/a : Y'/X' \to Y/X$ sont les diagrammes commutatifs
$$
\xymatrix{
Y' \ar[d] \ar[r]_b & Y \ar[d] \\
X' \ar[r]^a & X
}
$$
induisant un isomorphisme de $Y'$ avec un sous-schéma ouvert de
$X' \times_X Y$. La proposition signifie que, pour tout objet
$Y/X$ de $\mathcal{C}$, on dispose d'un isomorphisme
$c_{Y/X} : \det(\NL_{Y/X}) \to \omega_{Y/X}$
tel que $c_{Y/X}(\delta(\NL_{Y/X})) = \tau_{Y/X}$,
et que, pour tout morphisme $b/a : Y'/X' \to Y/X$ de $\mathcal{C}$, on a
$b^*c_{Y/X} = c_{Y'/X'}$ via les identifications
$b^*\det(\NL_{Y/X}) = \det(\NL_{Y'/X'})$ et
$b^*\omega_{Y/X} = \omega_{Y'/X'}$ décrites ci-dessus.

\medskip\noindent
Étant donné $Y/X$ dans $\mathcal{C}$ et $y \in Y$, on peut trouver
des ouverts affines $V \subset Y$ et $U \subset X$ avec
$y \in V$ et $f(V) \subset U$
tels qu'il existe un isomorphisme
$$
\det(\NL_{Y/X})|_V \longrightarrow \omega_{Y/X}|_V
$$
envoyant $\delta(\NL_{Y/X})|_V$ sur $\tau_{Y/X}|_V$. Cela résulte
du choix d'ouverts affines comme dans le
lemme \ref{lemma-syntomic-quasi-finite}, (5), de la description
locale affine de $\delta(\NL_{Y/X})$ dans la
remarque \ref{remark-local-description-delta}, et du
lemme \ref{lemma-different-quasi-finite-complete-intersection}.
Si l'annulateur de la section $\tau_{Y/X}$ est nul, ces
applications locales sont uniques et se recollent automatiquement. Ainsi, si l'annulateur
de $\tau_{Y/X}$ est nul, il existe un unique isomorphisme
$c_{Y/X} : \det(\NL_{Y/X}) \to \omega_{Y/X}$ tel que
$c_{Y/X}(\delta(\NL_{Y/X})) = \tau_{Y/X}$.
Si $b/a : Y'/X' \to Y/X$ est un morphisme de $\mathcal{C}$
et si l'annulateur de $\tau_{Y'/X'}$ est lui aussi nul,
alors $b^*c_{Y/X}$ est l'unique isomorphisme
$c_{Y'/X'} : \det(\NL_{Y'/X'}) \to \omega_{Y'/X'}$ tel que
$c_{Y'/X'}(\delta(\NL_{Y'/X'})) = \tau_{Y'/X'}$.
Cela découle formellement du fait que
$b^*\delta(\NL_{Y/X}) = \delta(\NL_{Y'/X'})$ et
$b^*\tau_{Y/X} = \tau_{Y'/X'}$.

\medskip\noindent
On peut résumer les résultats du paragraphe précédent comme suit.
Notons $\mathcal{C}_{nice} \subset \mathcal{C}$ la
sous-catégorie pleine des $Y/X$ tels que l'annulateur de
$\tau_{Y/X}$ soit nul. Nous avons alors résolu le problème
sur $\mathcal{C}_{nice}$. Pour $Y/X$ dans $\mathcal{C}_{nice}$,
nous continuons à noter $c_{Y/X}$ la solution que nous venons de trouver.

\medskip\noindent
Considérons des morphismes
$$
Y_1/X_1 \xleftarrow{b_1/a_1} Y/X \xrightarrow{b_2/a_2} Y_2/X_2
$$
dans $\mathcal{C}$ tels que $Y_1/X_1$ et $Y_2/X_2$ soient des objets
de $\mathcal{C}_{nice}$. {\bf Assertion.} $b_1^*c_{Y_1/X_1} = b_2^*c_{Y_2/X_2}$.
Nous montrerons d'abord que cette assertion entraîne la proposition,
puis nous la démontrerons.

\medskip\noindent
Soient $d, n \geq 1$ ; considérons le morphisme localement
quasi-fini syntomique $Y_{n, d} \to X_{n, d}$
construit dans l'exemple \ref{example-universal-quasi-finite-syntomic}.
Alors $Y_{n, d}$ est un schéma irréductible régulier et le
morphisme $Y_{n, d} \to X_{n, d}$ est localement quasi-fini syntomique
et étale au-dessus d'un ouvert dense, voir le
lemme \ref{lemma-universal-quasi-finite-syntomic-etale}.
Ainsi $\tau_{Y_{n, d}/X_{n, d}}$ est non nul, par exemple d'après le
lemme \ref{lemma-different-ramification}. Or une section non nulle
d'un module inversible sur un schéma irréductible régulier
a un annulateur nul. Ainsi
$Y_{n, d}/X_{n, d}$ est un objet de $\mathcal{C}_{nice}$.

\medskip\noindent
Soit $Y/X$ un objet quelconque de $\mathcal{C}$. Soit $y \in Y$.
D'après le lemme \ref{lemma-locally-comes-from-universal}, on peut trouver
$n, d \geq 1$ et des morphismes
$$
Y/X \leftarrow V/U \xrightarrow{b/a} Y_{n, d}/X_{n, d}
$$
de $\mathcal{C}$ tels que $V \subset Y$ et $U \subset X$ soient ouverts.
On peut donc prendre l'image réciproque du morphisme canonique $c_{Y_{n, d}/X_{n, d}}$
construit ci-dessus par $b$ sur $V$. L'assertion garantit que ces isomorphismes
locaux se recollent ! On obtient ainsi un isomorphisme global bien défini
$c_{Y/X} : \det(\NL_{Y/X}) \to \omega_{Y/X}$ tel que
$c_{Y/X}(\delta(\NL_{Y/X})) = \tau_{Y/X}$.
Si $b/a : Y'/X' \to Y/X$ est un morphisme de $\mathcal{C}$,
l'assertion entraîne également que l'application construite de façon analogue
$c_{Y'/X'}$ est l'image réciproque par $b$ de l'application construite localement
$c_{Y/X}$. Il reste donc à démontrer l'assertion.

\medskip\noindent
C'est ce que nous faisons dans la suite de la démonstration. On peut choisir un point
$y \in Y$ et montrer que les applications coïncident dans un voisinage ouvert de $y$.
On peut donc remplacer $Y_1$, $Y_2$ par des voisinages ouverts de
l'image de $y$ dans $Y_1$ et $Y_2$. On peut ainsi supposer que
$Y, X, Y_1, X_1, Y_2, X_2$ sont affines, disons les spectres
d'anneaux $B, A, B_1, A_1, B_2, A_2$. On a le diagramme
$$
\xymatrix{
B_1 \ar[r] & B & B_2 \ar[l] \\
A_1 \ar[u] \ar[r] & A \ar[u] & A_2 \ar[l] \ar[u]
}
$$
Par hypothèse, le spectre de $B$ est un ouvert affine
à la fois du spectre de $A \otimes_{A_1} B_1$ et de celui de $A \otimes_{A_2} B_2$.
En rétrécissant encore, on peut supposer qu'il existe des éléments
$g_i \in A \otimes_{A_i} B_i$ tels que nos applications donnent des isomorphismes
$(A \otimes_{A_i} B_i)_{g_i} = B$, voir
Propriétés, lemme \ref{properties-lemma-standard-open-two-affines}.
Soient $x_\alpha$, $y_\beta$ des collections assez grandes de
variables pour que l'on puisse choisir des surjections
$$
A'_1 = A_1[x_\alpha] \to A
\quad\text{et}\quad
A'_2 = A_2[y_\beta] \to A
$$
de $A_1$-algèbres et de $A_2$-algèbres. On peut alors choisir des relèvements
$h_i \in A'_i \otimes_{A_i} B_i$ de $g_i \in A \otimes_{A_i} B_i$,
et l'on considère le diagramme
$$
\xymatrix{
(A'_1 \otimes_{A_1} B_1)_{h_1} \ar[r] & B &
(A'_2 \otimes_{A_1} B_2)_{h_2} \ar[l] \\
A'_1 \ar[u] \ar[r] & A \ar[u] & A'_2 \ar[l] \ar[u]
}
$$
Par construction, les deux carrés sont cocartésiens. Considérons ensuite
l'homomorphisme d'anneaux
$$
A' = A'_1 \times_A A'_2 \longrightarrow
B' =
(A'_1 \otimes_{A_1} B_1)_{h_1} \times_B
(A'_2 \otimes_{A_1} B_2)_{h_2}
$$
D'après
Compléments d'algèbre, lemme \ref{more-algebra-lemma-module-over-fibre-product},
on a $A'_1 \otimes_{A'} B' = (A'_1 \otimes_{A_1} B_1)_{h_1}$
et $A'_2 \otimes_{A'} B' = (A'_2 \otimes_{A_2} B_2)_{h_2}$. En particulier,
les fibres du morphisme $Y' = \Spec(B') \to \Spec(A') = X'$ sont
des sous-schémas ouverts de changements de base des fibres des morphismes $Y_i \to X_i$,
donc finies et localement des intersections complètes
(voir le lemme \ref{lemma-syntomic-quasi-finite}).
D'après Compléments d'algèbre, lemme
\ref{more-algebra-lemma-properties-algebras-over-fibre-product},
l'homomorphisme d'anneaux $A' \to B'$ est plat et de présentation finie.
D'après le lemme \ref{lemma-syntomic-quasi-finite}, (3),
on conclut donc que $Y' \to X'$ est localement quasi-fini et syntomique.
On obtient ainsi le diagramme commutatif
$$
\xymatrix{
& & Y/X \ar[ld] \ar@/_2pc/[lldd]_{b_1/a_1} \ar[rd] \ar@/^2pc/[rrdd]^{b_2/a_2} \\
& Y'_1/X'_1 \ar[rd] \ar[ld] & & Y'_2/X'_2 \ar[ld] \ar[rd] \\
Y_1/X_1 & & Y'/X' & & Y_2/X_2
}
$$
où $Y'_i/X'_i$ correspond à $A'_i \to (A'_i \otimes_{A_i} B_i)_{h_i}$.

\medskip\noindent
Dans ce paragraphe, un passage à la limite remédie au fait que $A'$ et $A'_i$ ne sont pas
nécessairement noethériens ; le lecteur peut omettre ce paragraphe.
On peut trouver une $\mathbf{Z}$-sous-algèbre de type fini
$A'' \subset A'$ et un homomorphisme quasi-fini syntomique d'anneaux
$A'' \to B''$ tels que $B' = A' \otimes_{A''} B''$.
Cela résulte de Limites, lemmes
\ref{limits-lemma-descend-finite-presentation},
\ref{limits-lemma-descend-syntomic} et
\ref{limits-lemma-descend-quasi-finite}.
Les homomorphismes d'anneaux $A'' \to A'_1 = A_1[x_\alpha]$ se factorisent alors par une
sous-algèbre de polynômes à un nombre fini de variables $A_i[x_1, \ldots, x_n$, et l'homomorphisme
$A'' \to A'_2 = A_2[y_\beta]$ se factorise par $A_2[y_1, \ldots, y_m]$.
En agrandissant les ensembles de variables, on peut aussi supposer que
$h_1 \in A_1[x_1, \ldots, x_n] \otimes_{A_1} B_1$ et
$h_2 \in A_2[y_1, \ldots, y_m] \otimes_{A_2} B_2$.
En agrandissant encore les ensembles de variables, on peut supposer
que les applications $B'' \to (A'_i \otimes_{A_i} B_i)_{h_i}$
se factorisent par $(A_1[x_1, \ldots, x_n] \otimes_{A_1} B_1)_{h_1}$
et $(A_2[y_1, \ldots, y_m] \otimes_{A_2} B_2)_{h_2}$.
En remplaçant $Y'$, $X'$, $Y'_1$, $X'_1$, $Y'_2$, $X'_2$ par les spectres de
$B''$, $A''$, $(A_1[x_1, \ldots, x_n] \otimes_{A_1} B_1)_{h_1}$,
$A_1[x_1, \ldots, x_n]$, $(A_2[y_1, \ldots, y_m] \otimes_{A_2} B_2)_{h_2}$,
$A_2[y_1, \ldots, y_m]$, on obtient un diagramme comme ci-dessus
où tous les schémas sont noethériens.

\medskip\noindent
Remarquons que $Y'_i/X'_i$ est un objet de $\mathcal{C}_{nice}$, car il
s'obtient en prenant un sous-schéma ouvert du produit d'un espace affine
avec $Y_i/X_i$. En particulier, l'image réciproque de
$c_{Y_i/X_i}$ par $(b_i/a_i)$ est la même que l'image réciproque de
$c_{Y'_i/X'_i}$ par $Y/X \to Y'_i/X'_i$.
La démonstration serait alors achevée si $Y'/X'$ était un objet de $\mathcal{C}_{nice}$
(mais ce n'est probablement pas le cas). En effet, en prenant l'image réciproque de $c_{Y'/X'}$
le long des deux côtés du carré, on obtiendrait la conclusion voulue.
Pour contourner le fait que $Y'/X'$ n'appartient pas à $\mathcal{C}_{nice}$,
remarquons que les arguments ci-dessus montrent qu'après avoir éventuellement rétréci tous
les schémas $X, Y, X'_1, Y'_1, X'_2, Y'_2, X', Y'$, on peut trouver
$n, d \geq 1$ et prolonger le diagramme comme suit :
$$
\xymatrix{
& Y/X \ar[ld] \ar[rd] \\
Y'_1/X'_1 \ar[rd] & & Y'_2/X'_2 \ar[ld] \\
& Y'/X' \ar[d] \\
& Y_{n, d}/X_{n, d}
}
$$
On peut alors utiliser l'argument déjà donné en prenant
l'image réciproque de $c_{Y_{n, d}/X_{n, d}}$. La démonstration est achevée.
\end{proof}













\section{Une généralisation de la différente}
\label{section-different-generalization}

\noindent
Dans cette section, nous généralisons la définition \ref{definition-different}
pour englober tous les homomorphismes d'anneaux $A \to B$
pour lesquels la différente de Dedekind est définie et $1 \in \mathcal{L}_{B/A}$.
Expliquons d'abord la condition « $A \to B$ envoie les non-diviseurs de zéro sur
des non-diviseurs de zéro et induit un homomorphisme plat $Q(A) \to Q(A) \otimes_A B$ ».

\begin{lemma}
\label{lemma-explain-condition}
Soit $A \to B$ un homomorphisme d'anneaux noethériens. Considérons les conditions
\begin{enumerate}
\item les non-diviseurs de zéro de $A$ s'envoient sur des non-diviseurs de zéro de $B$ ;
\item (1) est satisfaite et $Q(A) \to Q(A) \otimes_A B$ est plat ;
\item $A \to B_\mathfrak q$ est plat pour tout
$\mathfrak q \in \text{Ass}(B)$ ;
\item (3) est satisfaite et $A \to B_\mathfrak q$ est plat pour tout $\mathfrak q$
au-dessus d'un élément de $\text{Ass}(A)$.
\end{enumerate}
On a alors les implications suivantes :
$$
\xymatrix{
(1) & (2) \ar@{=>}[l] \ar@{=>}[d] \\
(3) \ar@{=>}[u] & (4) \ar@{=>}[l]
}
$$
Si $A \to B$ vérifie la propriété de montée, (2) et (4) sont équivalentes.
\end{lemma}

\begin{proof}
Les implications horizontales du diagramme sont triviales.
Soit $S \subset A$ l'ensemble des non-diviseurs de zéro, de sorte que
$Q(A) = S^{-1}A$ et $Q(A) \otimes_A B = S^{-1}B$. Rappelons que
$S = A \setminus \bigcup_{\mathfrak p \in \text{Ass}(A)} \mathfrak p$
d'après Algèbre, lemme \ref{algebra-lemma-ass-zero-divisors}.
Soit $\mathfrak q \subset B$ un idéal premier au-dessus de $\mathfrak p \subset A$.

\medskip\noindent
Supposons (2). Si $\mathfrak q \in \text{Ass}(B)$, alors
$\mathfrak q$ est constitué de diviseurs de zéro ; (1) entraîne donc
qu'il en est de même de $\mathfrak p$. Ainsi
$\mathfrak p$ correspond à un idéal premier de $S^{-1}A$.
Par conséquent, $A \to B_\mathfrak q$ est plat d'après l'hypothèse (2).
Si $\mathfrak q$ est au-dessus d'un idéal premier associé $\mathfrak p$
de $A$, alors on a certainement $\mathfrak p \in \Spec(S^{-1}A)$, et le
même argument s'applique.

\medskip\noindent
Supposons (3). Soit $f \in A$ un non-diviseur de zéro. Si $f$ était diviseur de zéro
sur $B$, alors $f$ appartiendrait à un idéal premier associé $\mathfrak q$
de $B$. Puisque $A \to B_\mathfrak q$ est plat par hypothèse, on en déduit que
$\mathfrak p$ est un idéal premier associé de $A$, d'après
Algèbre, lemme \ref{algebra-lemma-bourbaki}. Cela entraînerait que
$f$ est diviseur de zéro sur $A$, contradiction.

\medskip\noindent
Supposons (4) et la propriété de montée pour $A \to B$. Nous savons déjà que (1) est satisfaite.
Si $\mathfrak q$ correspond à un idéal premier de $S^{-1}B$, alors $\mathfrak p$
est contenu dans un idéal premier associé $\mathfrak p'$ de $A$. Par la propriété de montée,
il existe un idéal premier $\mathfrak q'$ contenant $\mathfrak q$ et situé
au-dessus de $\mathfrak p$. Alors $A \to B_{\mathfrak q'}$ est plat d'après
(4). Donc $A \to B_{\mathfrak q}$ est plat, en tant que localisation.
Ainsi $A \to S^{-1}B$ est plat, de même que $S^{-1}A \to S^{-1}B$, voir
Algèbre, lemme \ref{algebra-lemma-flat-localization}.
\end{proof}

\begin{remark}
\label{remark-different-generalization}
On peut généraliser la définition \ref{definition-different}.
Supposons que $f : Y \to X$ soit un morphisme quasi-fini de schémas noethériens
possédant les propriétés suivantes :
\begin{enumerate}
\item l'ouvert $V \subset Y$ où $f$ est plat contient
$\text{Ass}(\mathcal{O}_Y)$ et $f^{-1}(\text{Ass}(\mathcal{O}_X))$ ;
\item l'élément trace $\tau_{V/X}$ provient d'une section
$\tau \in \Gamma(Y, \omega_{Y/X})$.
\end{enumerate}
La condition (1) entraîne que $V$ contient les points associés de
$\omega_{Y/X}$, d'après le lemme \ref{lemma-dualizing-associated-primes}.
En particulier, $\tau$ est unique s'il existe
(Diviseurs, lemme \ref{divisors-lemma-restriction-injective-open-contains-ass}).
Étant donné $\tau$, on peut définir la différente $\mathfrak{D}_f$ comme l'annulateur de
$\Coker(\tau : \mathcal{O}_Y \to \omega_{Y/X})$. Elle coïncide avec la
différente de Dedekind dans de nombreux cas (lemme \ref{lemma-agree-dedekind}).
Cependant, pour les homomorphismes non plats entre anneaux non normaux, cette généralisation
ne mesure plus la ramification du morphisme, voir
l'exemple \ref{example-no-different}.
\end{remark}

\begin{lemma}
\label{lemma-agree-dedekind}
Supposons que la différente de Dedekind de $A \to B$ soit définie.
Posons $X = \Spec(A)$ et $Y = \Spec(B)$. La généralisation de la
remarque \ref{remark-different-generalization}
s'applique au morphisme $f : Y \to X$ si et seulement si
$1 \in \mathcal{L}_{B/A}$ (par exemple si $A$ est normal, voir le
lemme \ref{lemma-dedekind-different-ideal}).
Dans ce cas, $\mathfrak{D}_{B/A}$ est un idéal de $B$ et l'on a
$$
\mathfrak{D}_f = \widetilde{\mathfrak{D}_{B/A}}
$$
comme faisceaux cohérents d'idéaux sur $Y$.
\end{lemma}

\begin{proof}
La différente de Dedekind de $A \to B$ étant définie, on peut appliquer le
lemme \ref{lemma-explain-condition} pour voir que
$Y \to X$ vérifie la condition (1) de la
remarque \ref{remark-different-generalization}.
Rappelons qu'il existe un isomorphisme canonique
$c : \mathcal{L}_{B/A} \to \omega_{B/A}$, voir le
lemme \ref{lemma-dedekind-complementary-module}.
Soient $K = Q(A)$ et $L = K \otimes_A B$ comme ci-dessus.
Par construction, l'application $c$ s'insère dans un diagramme commutatif
$$
\xymatrix{
\mathcal{L}_{B/A} \ar[r] \ar[d]_c & L \ar[d] \\
\omega_{B/A} \ar[r] & \Hom_K(L, K)
}
$$
où la flèche verticale de droite envoie $x \in L$ sur l'application
$y \mapsto \text{Trace}_{L/K}(xy)$, et la flèche horizontale
inférieure est l'application de changement de base (\ref{equation-bc-dualizing}) pour $\omega_{B/A}$.
On peut factoriser l'application horizontale inférieure en
$$
\omega_{B/A} = \Gamma(Y, \omega_{Y/X})
\to \Gamma(V, \omega_{V/X}) \to \Hom_K(L, K)
$$
Puisque tous les points associés de $\omega_{V/X}$
s'envoient sur des idéaux premiers associés de $A$
(lemme \ref{lemma-dualizing-associated-primes}),
la seconde application est injective.
L'élément $\tau_{V/X}$ s'envoie sur $\text{Trace}_{L/K}$ dans
$\Hom_K(L, K)$ par la définition même des éléments trace
(définition \ref{definition-trace-element}).
Ainsi un $\tau$ comme dans la condition (2) de la
remarque \ref{remark-different-generalization}
existe si et seulement si $1 \in \mathcal{L}_{B/A}$, et l'on a alors
$\tau = c(1)$. Dans ce cas, d'après le lemme \ref{lemma-dedekind-different-ideal},
on voit que $\mathfrak{D}_{B/A} \subset B$.
Enfin, l'égalité de $\mathfrak{D}_f$ et de $\mathfrak{D}_{B/A}$
découle immédiatement des définitions et du fait que $\tau = c(1)$, établi ci-dessus.
\end{proof}

\begin{example}
\label{example-no-different}
Soit $k$ un corps. Soient $A = k[x, y]/(xy)$ et $B = k[u, v]/(uv)$, et soit
$A \to B$ donné par $x \mapsto u^n$ et $y \mapsto v^m$, pour des
$n, m \in \mathbf{N}$ premiers à la caractéristique de $k$. Alors
$A_{x + y} \to B_{x + y}$ est (fini) étale, et nous sommes donc dans une situation
où la différente de Dedekind est définie. Un calcul donne
$$
\text{Trace}_{L/K}(1) = (nx + my)/(x + y),\quad
\text{Trace}_{L/K}(u^i) = 0,\quad \text{Trace}_{L/K}(v^j) = 0
$$
pour $1 \leq i < n$ et $1 \leq j < m$. On en déduit que
$1 \in \mathcal{L}_{B/A}$ si et seulement si $n = m$. De plus, un
calcul montre que, si $n = m$, alors $\mathcal{L}_{B/A} = B$,
et la différente de Dedekind est également $B$. Autrement dit,
la différente de la remarque \ref{remark-different-generalization}
est définie pour $\Spec(B) \to \Spec(A)$
si et seulement si $n = m$, et, dans ce cas, elle est
l'idéal unité. On voit donc que, dans les cas non plats, la non-annulation
de la différente ne garantit pas que le morphisme soit étale ou non ramifié.
\end{example}







\section{Comparaison avec la théorie de la dualité}
\label{section-comparison}

\noindent
Dans cette section, nous comparons les constructions algébriques élémentaires
ci-dessus aux constructions du chapitre sur la théorie de la dualité
des schémas.

\begin{lemma}
\label{lemma-compare-dualizing}
Soit $f : Y \to X$ un morphisme quasi-fini séparé de schémas noethériens.
Pour tout couple d'ouverts affines $\Spec(B) = V \subset Y$,
$\Spec(A) = U \subset X$ avec $f(V) \subset U$, il existe un isomorphisme
$$
H^0(V, f^!\mathcal{O}_X) = \omega_{B/A}
$$
où $f^!$ est comme dans
Dualité des schémas, section \ref{duality-section-upper-shriek}.
Ces isomorphismes sont compatibles aux applications de restriction et définissent un
isomorphisme canonique $H^0(f^!\mathcal{O}_X) = \omega_{Y/X}$, avec
$\omega_{Y/X}$ comme dans la remarque \ref{remark-relative-dualizing-for-quasi-finite}.
De même, si $f : Y \to X$ est un morphisme quasi-fini de schémas de
type fini sur une base noethérienne $S$ munie d'un complexe dualisant
$\omega_S^\bullet$, alors $H^0(f_{new}^!\mathcal{O}_X) = \omega_{Y/X}$.
\end{lemma}

\begin{proof}
Le théorème principal de Zariski permet de choisir une factorisation $f = f' \circ j$,
où $j : Y \to Y'$ est une immersion ouverte et $f' : Y' \to X$ un morphisme
fini, voir Compléments sur les morphismes, lemme
\ref{more-morphisms-lemma-quasi-finite-separated-pass-through-finite}.
Par la construction de
Dualité des schémas, lemme \ref{duality-lemma-shriek-well-defined}, on a
$f^! = j^* \circ a'$, où
$a' : D_\QCoh(\mathcal{O}_X) \to D_\QCoh(\mathcal{O}_{Y'})$
est l'adjoint à droite de $Rf'_*$ donné par
Dualité des schémas, lemme \ref{duality-lemma-twisted-inverse-image}.
D'après Dualité des schémas, lemme \ref{duality-lemma-finite-twisted},
on voit que
$\Phi(a'(\mathcal{O}_X)) = R\SheafHom(f'_*\mathcal{O}_{Y'}, \mathcal{O}_X)$ dans
$D_\QCoh^+(f'_*\mathcal{O}_{Y'})$. En particulier, $a'(\mathcal{O}_X)$ a des
faisceaux de cohomologie nuls en degrés $< 0$. Le faisceau de cohomologie de degré zéro
est déterminé par l'isomorphisme
$$
f'_*H^0(a'(\mathcal{O}_X)) =
\SheafHom_{\mathcal{O}_X}(f'_*\mathcal{O}_{Y'}, \mathcal{O}_X)
$$
de $f'_*\mathcal{O}_{Y'}$-modules via l'équivalence de
Morphismes, lemme \ref{morphisms-lemma-affine-equivalence-modules}.\newline
En écrivant $(f')^{-1}U = V' = \Spec(B')$, on obtient
$$
H^0(V', a'(\mathcal{O}_X)) = \Hom_A(B', A).
$$
Puisque le faisceau de cohomologie de degré zéro de $a'(\mathcal{O}_X)$
est un module quasi-cohérent, on trouve que
sa restriction à $V$ est donnée par
$\omega_{B/A} = \Hom_A(B', A) \otimes_{B'} B$, comme voulu.

\medskip\noindent
L'assertion sur les applications de restriction signifie que celles
du $\mathcal{O}_{Y'}$-module quasi-cohérent $H^0(a'(\mathcal{O}_X))$
pour les ouverts de $Y'$ coïncident avec les applications définies dans le
lemme \ref{lemma-localize-dualizing}
pour les modules $\omega_{B/A}$ via les isomorphismes ci-dessus.
C'est clair.

\medskip\noindent
Soit $f : Y \to X$ un morphisme quasi-fini de schémas de type fini
sur une base noethérienne $S$ munie d'un complexe dualisant $\omega_S^\bullet$.
Considérons des ouverts $V \subset Y$ et $U \subset X$ avec $f(V) \subset U$,
et $V$ et $U$ séparés sur $S$. Notons $f|_V : V \to U$ la restriction
de $f$. D'après la discussion ci-dessus et
Dualité des schémas, lemme \ref{duality-lemma-duality-bootstrap},
on a des isomorphismes canoniques
$$
H^0(f_{new}^!\mathcal{O}_X)|_V = H^0((f|_V)^!\mathcal{O}_U) = \omega_{V/U} =
\omega_{Y/X}|_V
$$
Nous omettons la vérification que ces isomorphismes se recollent en un isomorphisme
global $H^0(f_{new}^!\mathcal{O}_X) \to \omega_{Y/X}$.
\end{proof}

\begin{lemma}
\label{lemma-compare-trace}
Soit $f : Y \to X$ un morphisme fini plat de schémas noethériens.
L'application
$$
\text{Trace}_f : f_*\mathcal{O}_Y \longrightarrow \mathcal{O}_X
$$
de la section \ref{section-discriminant}
correspond à une application $\mathcal{O}_Y \to f^!\mathcal{O}_X$ (voir la démonstration).
Notons $\tau_{Y/X} \in H^0(Y, f^!\mathcal{O}_X)$ l'image de $1$.
Via l'isomorphisme $H^0(f^!\mathcal{O}_X) = \omega_{X/Y}$ du
lemme \ref{lemma-compare-dualizing},
cela coïncide avec la construction de la
remarque \ref{remark-relative-dualizing-for-flat-quasi-finite}.
\end{lemma}

\begin{proof}
Le foncteur $f^!$ est défini dans
Dualité des schémas, section \ref{duality-section-upper-shriek}.
Puisque $f$ est fini (donc propre), $f^!$ est donné par
l'adjoint à droite de l'image directe par $f$. Dans
Dualité des schémas, section \ref{duality-section-duality-finite},
nous avons explicité cet adjoint. En particulier,
l'objet $f^!\mathcal{O}_X$ consiste en un seul
faisceau de cohomologie placé en degré $0$, et l'on a pour ce faisceau
$$
f_*f^!\mathcal{O}_X =
\SheafHom_{\mathcal{O}_X}(f_*\mathcal{O}_Y, \mathcal{O}_X)
$$
Pour le voir, on utilise aussi le fait que $f_*\mathcal{O}_Y$ est fini localement libre,
car $f$ est un morphisme fini plat de schémas noethériens,
et donc que tous les faisceaux Ext supérieurs sont nuls. Certains détails sont omis.
Finalement,
$$
\text{Trace}_f \in
\Hom_{\mathcal{O}_X}(f_*\mathcal{O}_Y, \mathcal{O}_X) =
\Gamma(X, f_*f^!\mathcal{O}_X) =
\Gamma(Y, f^!\mathcal{O}_X)
$$
D'autre part, on a $f^!\mathcal{O}_X = \omega_{Y/X}$
par l'identification du lemme \ref{lemma-compare-dualizing}.
On dispose donc maintenant de deux éléments, à savoir $\text{Trace}_f$
et $\tau_{Y/X}$ de la
remarque \ref{remark-relative-dualizing-for-flat-quasi-finite}, dans
$$
\Gamma(Y, f^!\mathcal{O}_X) = \Gamma(Y, \omega_{Y/X})
$$
et le lemme affirme qu'ils sont égaux.

\medskip\noindent
Soit $U = \Spec(A) \subset X$ un ouvert affine d'image réciproque
$V = \Spec(B) \subset Y$. Puisque $f$ est fini,
$A \to B$ est fini et l'on a donc $\omega_{Y/X}(V) = \Hom_A(B,A)$
par construction ; cet isomorphisme coïncide avec l'identification
de $f_*f^!\mathcal{O}_Y$ avec
$\SheafHom_{\mathcal{O}_X}(f_*\mathcal{O}_Y, \mathcal{O}_X)$ décrite
ci-dessus. L'égalité de $\text{Trace}_f$ et de $\tau_{Y/X}$
découle donc du fait que $\tau_{B/A} = \text{Trace}_{B/A}$
d'après le lemme \ref{lemma-finite-flat-trace}.
\end{proof}








\section{Morphismes quasi-finis de Gorenstein}
\label{section-gorenstein-lci}

\noindent
Cette section traite des morphismes quasi-finis de Gorenstein.

\begin{lemma}
\label{lemma-gorenstein-quasi-finite}
Soit $f : Y \to X$ un morphisme quasi-fini de schémas noethériens.
Les conditions suivantes sont équivalentes :
\begin{enumerate}
\item $f$ est de Gorenstein ;
\item $f$ est plat et les fibres de $f$ sont de Gorenstein ;
\item $f$ est plat et $\omega_{Y/X}$ est inversible
(remarque \ref{remark-relative-dualizing-for-quasi-finite}) ;
\item pour tout $y \in Y$, il existe des ouverts affines
$y \in V = \Spec(B) \subset Y$, $U = \Spec(A) \subset X$
avec $f(V) \subset U$ tels que $A \to B$ soit plat
et $\omega_{B/A}$ soit un $B$-module inversible.
\end{enumerate}
\end{lemma}

\begin{proof}
Les conditions (1) et (2) sont équivalentes par définition. Les conditions (3) et (4)
sont équivalentes par la construction de $\omega_{Y/X}$ dans la
remarque \ref{remark-relative-dualizing-for-quasi-finite}.
Il faut donc montrer que (1)-(2) équivaut à (3)-(4).

\medskip\noindent
Première démonstration. En travaillant localement sur les ouverts affines, on peut supposer $f$ séparé
et appliquer le lemme \ref{lemma-compare-dualizing} pour voir que
$\omega_{Y/X}$ est le faisceau de cohomologie de degré zéro de $f^!\mathcal{O}_X$.
Sous chacune des deux hypothèses, $f$ est plat et quasi-fini ;
$f^!\mathcal{O}_X$ est donc isomorphe à $\omega_{Y/X}[0]$, voir
Dualité des schémas, lemme \ref{duality-lemma-flat-quasi-finite-shriek}. Ainsi
l'équivalence découle de
Dualité des schémas, lemme
\ref{duality-lemma-affine-flat-Noetherian-gorenstein}.

\medskip\noindent
Deuxième démonstration. D'après le lemme \ref{lemma-characterize-invertible},
il suffit d'établir l'équivalence de
(2) et (3) lorsque $X$ est le spectre d'un corps $k$.
Alors $Y = \Spec(B)$, où $B$ est une $k$-algèbre finie.
Dans ce cas, $\omega_{B/A} = \omega_{B/k} = \Hom_k(B, k)$,
placé en degré $0$, est un complexe dualisant de $B$, voir
Complexes dualisants, lemme \ref{dualizing-lemma-dualizing-finite}.
L'équivalence découle donc de
Complexes dualisants, lemme \ref{dualizing-lemma-gorenstein}.
\end{proof}

\begin{remark}
\label{remark-collect-results-qf-gorenstein}
Soit $f : Y \to X$ un morphisme quasi-fini de Gorenstein de schémas noethériens.
Soient $\mathfrak D_f \subset \mathcal{O}_Y$ la différente et
$R \subset Y$ le sous-schéma fermé défini par $\mathfrak D_f$.
On a alors
\begin{enumerate}
\item $\mathfrak D_f$ est un idéal localement principal ;
\item $R$ est un sous-schéma fermé localement principal ;
\item $\mathfrak D_f$ coïncide, localement sur les ouverts affines, avec la différente de Noether ;
\item la formation de $R$ commute au changement de base ;
\item si $f$ est fini, la norme de $R$ est le discriminant de $f$ ;
\item si $f$ est étale aux points associés de $Y$, alors
$R$ est un diviseur de Cartier effectif et $\omega_{Y/X} = \mathcal{O}_Y(R)$.
\end{enumerate}
Cela résulte des lemmes \ref{lemma-flat-gorenstein-agree-noether},
\ref{lemma-base-change-different} et
\ref{lemma-norm-different-is-discriminant}.
\end{remark}

\begin{remark}
\label{remark-collect-results-qf-gorenstein-two}
Soit $S$ un schéma noethérien muni d'un complexe dualisant
$\omega_S^\bullet$. Soit $f : Y \to X$ un morphisme quasi-fini de Gorenstein
de schémas compactifiables sur $S$. Supposons de plus
$Y$ et $X$ de Cohen-Macaulay et $f$ étale aux points génériques
de $Y$. On peut alors combiner
Dualité des schémas, remarque
\ref{duality-remark-CM-morphism-compare-dualizing} et la
remarque \ref{remark-collect-results-qf-gorenstein}
pour obtenir un isomorphisme canonique
$$
\omega_Y = f^*\omega_X \otimes_{\mathcal{O}_Y} \omega_{Y/X} =
f^*\omega_X \otimes_{\mathcal{O}_Y} \mathcal{O}_Y(R)
$$
de $\mathcal{O}_Y$-modules. Si $f$ est en outre fini,
l'isomorphisme $\mathcal{O}_Y(R) = \omega_{Y/X}$ provient
de la section globale $\tau_{Y/X} \in H^0(Y, \omega_{Y/X})$
qui correspond par dualité à l'application
$\text{Trace}_f : f_*\mathcal{O}_Y \to \mathcal{O}_X$, voir le
lemme \ref{lemma-compare-trace}.
\end{remark}






\begin{multicols}{2}[\section{Autres chapitres}]
\noindent
Préliminaires
\begin{enumerate}
\item \hyperref[introduction-section-phantom]{Introduction}
\item \hyperref[conventions-section-phantom]{Conventions}
\item \hyperref[sets-section-phantom]{Théorie des ensembles}
\item \hyperref[categories-section-phantom]{Catégories}
\item \hyperref[topology-section-phantom]{Topologie}
\item \hyperref[sheaves-section-phantom]{Faisceaux sur les espaces}
\item \hyperref[sites-section-phantom]{Sites et faisceaux}
\item \hyperref[stacks-section-phantom]{Champs}
\item \hyperref[fields-section-phantom]{Corps}
\item \hyperref[algebra-section-phantom]{Algèbre commutative}
\item \hyperref[brauer-section-phantom]{Groupes de Brauer}
\item \hyperref[homology-section-phantom]{Algèbre homologique}
\item \hyperref[derived-section-phantom]{Catégories dérivées}
\item \hyperref[simplicial-section-phantom]{Méthodes simpliciales}
\item \hyperref[more-algebra-section-phantom]{Compléments d'algèbre}
\item \hyperref[smoothing-section-phantom]{Lissage des morphismes d'anneaux}
\item \hyperref[modules-section-phantom]{Faisceaux de modules}
\item \hyperref[sites-modules-section-phantom]{Modules sur les sites}
\item \hyperref[injectives-section-phantom]{Objets injectifs}
\item \hyperref[cohomology-section-phantom]{Cohomologie des faisceaux}
\item \hyperref[sites-cohomology-section-phantom]{Cohomologie sur les sites}
\item \hyperref[dga-section-phantom]{Algèbre différentielle graduée}
\item \hyperref[dpa-section-phantom]{Algèbre à puissances divisées}
\item \hyperref[sdga-section-phantom]{Faisceaux différentiels gradués}
\item \hyperref[hypercovering-section-phantom]{Hyperrecouvrements}
\end{enumerate}
Schémas
\begin{enumerate}
\setcounter{enumi}{25}
\item \hyperref[schemes-section-phantom]{Schémas}
\item \hyperref[constructions-section-phantom]{Constructions de schémas}
\item \hyperref[properties-section-phantom]{Propriétés des schémas}
\item \hyperref[morphisms-section-phantom]{Morphismes de schémas}
\item \hyperref[coherent-section-phantom]{Cohomologie des schémas}
\item \hyperref[divisors-section-phantom]{Diviseurs}
\item \hyperref[limits-section-phantom]{Limites de schémas}
\item \hyperref[varieties-section-phantom]{Variétés}
\item \hyperref[topologies-section-phantom]{Topologies sur les schémas}
\item \hyperref[descent-section-phantom]{Descente}
\item \hyperref[perfect-section-phantom]{Catégories dérivées de schémas}
\item \hyperref[more-morphisms-section-phantom]{Compléments sur les morphismes}
\item \hyperref[flat-section-phantom]{Compléments sur la platitude}
\item \hyperref[groupoids-section-phantom]{Schémas en groupoïdes}
\item \hyperref[more-groupoids-section-phantom]{Compléments sur les schémas en groupoïdes}
\item \hyperref[etale-section-phantom]{Morphismes étales de schémas}
\end{enumerate}
Thèmes de théorie des schémas
\begin{enumerate}
\setcounter{enumi}{41}
\item \hyperref[chow-section-phantom]{Homologie de Chow}
\item \hyperref[intersection-section-phantom]{Théorie de l'intersection}
\item \hyperref[pic-section-phantom]{Schémas de Picard des courbes}
\item \hyperref[weil-section-phantom]{Théories cohomologiques de Weil}
\item \hyperref[adequate-section-phantom]{Modules adéquats}
\item \hyperref[dualizing-section-phantom]{Complexes dualisants}
\item \hyperref[duality-section-phantom]{Dualité pour les schémas}
\item \hyperref[discriminant-section-phantom]{Discriminants et différentes}
\item \hyperref[derham-section-phantom]{Cohomologie de de Rham}
\item \hyperref[local-cohomology-section-phantom]{Cohomologie locale}
\item \hyperref[algebraization-section-phantom]{Géométrie algébrique et formelle}
\item \hyperref[curves-section-phantom]{Courbes algébriques}
\item \hyperref[resolve-section-phantom]{Résolution des surfaces}
\item \hyperref[models-section-phantom]{Réduction semi-stable}
\item \hyperref[functors-section-phantom]{Foncteurs et morphismes}
\item \hyperref[equiv-section-phantom]{Catégories dérivées de variétés}
\item \hyperref[pione-section-phantom]{Groupes fondamentaux de schémas}
\item \hyperref[etale-cohomology-section-phantom]{Cohomologie étale}
\item \hyperref[crystalline-section-phantom]{Cohomologie cristalline}
\item \hyperref[proetale-section-phantom]{Cohomologie pro-étale}
\item \hyperref[relative-cycles-section-phantom]{Cycles relatifs}
\item \hyperref[more-etale-section-phantom]{Compléments de cohomologie étale}
\item \hyperref[trace-section-phantom]{La formule des traces}
\end{enumerate}
Espaces algébriques
\begin{enumerate}
\setcounter{enumi}{64}
\item \hyperref[spaces-section-phantom]{Espaces algébriques}
\item \hyperref[spaces-properties-section-phantom]{Propriétés des espaces algébriques}
\item \hyperref[spaces-morphisms-section-phantom]{Morphismes d'espaces algébriques}
\item \hyperref[decent-spaces-section-phantom]{Espaces algébriques décents}
\item \hyperref[spaces-cohomology-section-phantom]{Cohomologie des espaces algébriques}
\item \hyperref[spaces-limits-section-phantom]{Limites d'espaces algébriques}
\item \hyperref[spaces-divisors-section-phantom]{Diviseurs sur les espaces algébriques}
\item \hyperref[spaces-over-fields-section-phantom]{Espaces algébriques sur des corps}
\item \hyperref[spaces-topologies-section-phantom]{Topologies sur les espaces algébriques}
\item \hyperref[spaces-descent-section-phantom]{Descente et espaces algébriques}
\item \hyperref[spaces-perfect-section-phantom]{Catégories dérivées d'espaces}
\item \hyperref[spaces-more-morphisms-section-phantom]{Compléments sur les morphismes d'espaces}
\item \hyperref[spaces-flat-section-phantom]{Platitude sur les espaces algébriques}
\item \hyperref[spaces-groupoids-section-phantom]{Groupoïdes dans les espaces algébriques}
\item \hyperref[spaces-more-groupoids-section-phantom]{Compléments sur les groupoïdes dans les espaces}
\item \hyperref[bootstrap-section-phantom]{Amorçage}
\item \hyperref[spaces-pushouts-section-phantom]{Sommes amalgamées d'espaces algébriques}
\end{enumerate}
Thèmes de géométrie
\begin{enumerate}
\setcounter{enumi}{81}
\item \hyperref[spaces-chow-section-phantom]{Groupes de Chow des espaces}
\item \hyperref[groupoids-quotients-section-phantom]{Quotients de groupoïdes}
\item \hyperref[spaces-more-cohomology-section-phantom]{Compléments sur la cohomologie des espaces}
\item \hyperref[spaces-simplicial-section-phantom]{Espaces simpliciaux}
\item \hyperref[spaces-duality-section-phantom]{Dualité pour les espaces}
\item \hyperref[formal-spaces-section-phantom]{Espaces algébriques formels}
\item \hyperref[restricted-section-phantom]{Algébrisation des espaces formels}
\item \hyperref[spaces-resolve-section-phantom]{Résolution des surfaces revisitée}
\end{enumerate}
Théorie des déformations
\begin{enumerate}
\setcounter{enumi}{89}
\item \hyperref[formal-defos-section-phantom]{Théorie formelle des déformations}
\item \hyperref[defos-section-phantom]{Théorie des déformations}
\item \hyperref[cotangent-section-phantom]{Le complexe cotangent}
\item \hyperref[examples-defos-section-phantom]{Problèmes de déformation}
\end{enumerate}
Champs algébriques
\begin{enumerate}
\setcounter{enumi}{93}
\item \hyperref[algebraic-section-phantom]{Champs algébriques}
\item \hyperref[examples-stacks-section-phantom]{Exemples de champs}
\item \hyperref[stacks-sheaves-section-phantom]{Faisceaux sur les champs algébriques}
\item \hyperref[criteria-section-phantom]{Critères de représentabilité}
\item \hyperref[artin-section-phantom]{Axiomes d'Artin}
\item \hyperref[quot-section-phantom]{Espaces de Quot et de Hilbert}
\item \hyperref[stacks-properties-section-phantom]{Propriétés des champs algébriques}
\item \hyperref[stacks-morphisms-section-phantom]{Morphismes de champs algébriques}
\item \hyperref[stacks-limits-section-phantom]{Limites de champs algébriques}
\item \hyperref[stacks-cohomology-section-phantom]{Cohomologie des champs algébriques}
\item \hyperref[stacks-perfect-section-phantom]{Catégories dérivées de champs}
\item \hyperref[stacks-introduction-section-phantom]{Introduction aux champs algébriques}
\item \hyperref[stacks-more-morphisms-section-phantom]{Compléments sur les morphismes de champs}
\item \hyperref[stacks-geometry-section-phantom]{La géométrie des champs}
\end{enumerate}
Thèmes de théorie des espaces de modules
\begin{enumerate}
\setcounter{enumi}{107}
\item \hyperref[moduli-section-phantom]{Champs de modules}
\item \hyperref[moduli-curves-section-phantom]{Espaces de modules de courbes}
\end{enumerate}
Divers
\begin{enumerate}
\setcounter{enumi}{109}
\item \hyperref[examples-section-phantom]{Exemples}
\item \hyperref[exercises-section-phantom]{Exercices}
\item \hyperref[guide-section-phantom]{Guide bibliographique}
\item \hyperref[desirables-section-phantom]{Desiderata}
\item \hyperref[coding-section-phantom]{Style de programmation}
\item \hyperref[obsolete-section-phantom]{Obsolète}
\item \hyperref[fdl-section-phantom]{Licence de documentation libre GNU}
\item \hyperref[index-section-phantom]{Index généré automatiquement}
\end{enumerate}
\end{multicols}


\bibliography{my}
\bibliographystyle{amsalpha}

\end{document}
